\documentclass[12pt,a4paper]{article}
\usepackage[spanish,es-tabla,es-noquoting]{babel}
\usepackage[utf8]{inputenc}
\usepackage[T1]{fontenc}
\usepackage{amsfonts}

\renewcommand{\rmdefault}{ptm}
\usepackage{mathptmx}
\DeclareSymbolFont{Xlargesymbols}{OMX}{cmex}{m}{n}
\DeclareMathSymbol{\Xsum}{\mathop}{Xlargesymbols}{80}

\DeclareSymbolFont{newtxletters}{OML}{ntxmi}{m}{it}
\SetSymbolFont{newtxletters}{bold}{OML}{ntxmi}{b}{it}
\DeclareMathSymbol{\mu}{\mathord}{newtxletters}{22}

\DeclareSymbolFont{Symbols}{OMS}{zplm}{m}{n}% Palatino
\DeclareMathSymbol{\infty}{\mathord}{Symbols}{"31}

\usepackage{lscape}
\usepackage{amsthm,tabls}
\usepackage{midpage}
\usepackage{subfigure}
\usepackage{changepage}
\usepackage{longtable}
\usepackage{graphicx,xspace}
\usepackage[a4paper,left=3.54cm,right=2.54cm,top=2.54cm,bottom=2.54cm,centering]{geometry}
\usepackage{makecell}
\newcolumntype{M}[1]{>{\centering\arraybackslash}m{#1}}
\usepackage{multirow}
%\usepackage{showframe}
\usepackage{bigstrut}
\renewcommand\theadfont{\normalsize}
\usepackage{csquotes}
\usepackage{titlesec}
\usepackage{enumitem}
\usepackage{booktabs}
\usepackage{setspace}
\usepackage{multirow}
\usepackage{rotating}

\usepackage{fancyhdr}
\pagestyle{fancyplain}
\fancyhf{}
\fancyhead[R]{\thepage}
\renewcommand{\headrulewidth}{0pt}

\usepackage[apaciteclassic, numberedbib, nosectionbib, tocbib]{apacite}
\usepackage{usebib}
\bibinput{une}
%\printbibliography[heading=none] with apacite
\makeatletter
\let\st@rtbibsection\@bibnewpage
\let\st@rtbibchapter\@bibnewpage
\makeatother
%\printbibliography[heading=none] with apacite

\usepackage[tablewithout,belowskip=-15pt,aboveskip=0pt]{caption}
\renewcommand{\thefigure}{\arabic{figure}}
\renewcommand{\thetable}{\arabic{table}}
\usepackage{url}
\newcommand{\ce}{\centering}
\def\UrlBreaks{\do\.\do\@\do\\\do\/\do\!\do\%\do\7\do\D\do\_\do\\do\;\do\>\do\]%
  \do\)\do\,\do\?\do\'\do+\do\=\do\#}%
\def\UrlBigBreaks{\do\:\do@url@hyp}%
\def\UrlBreaks{\do\/\do-}
%\DefineBibliographyStrings{spanish}{andothers={et~al\adddot}}

\renewcommand{\thesection}{\Roman{section}}
\renewcommand{\thesubsection}{\arabic{section}.\arabic{subsection}}

\titleformat*{\section}{\normalsize\bfseries}
%\setcounter{secnumdepth}{5}
\newcommand\Tstrut{\rule{0pt}{2.35ex}}
\newcommand\Bstrut{\rule[-1.3ex]{0pt}{1.3ex}}
\newcommand{\TBstrut}{\Tstrut\Bstrut}

\renewcommand{\qedsymbol}{$\blacksquare$}

\newcommand{\ww}{0.015em}
\titlespacing*{\chapter}
{0pt} {5em plus \ww minus \ww} {0em plus \ww minus \ww}
\titlespacing*{\section}
{0pt} {0em plus \ww minus \ww}{0em plus \ww minus \ww}
\titlespacing*{\subsection}
{0pt} {0em plus \ww minus \ww}{0em plus \ww minus \ww} 
\titlespacing*{\subsubsection}
{0pt} {0em plus \ww minus \ww}{0em plus \ww minus \ww}
\titlespacing*{\paragraph} {0pt}{1.25ex plus 1ex minus .2ex}{2em}
\titlespacing*{\subparagraph} {\parindent}{3.25ex plus 1ex minus .2ex}{1em}
\titleformat*{\section}{\normalsize\bfseries}
\titleformat*{\subsection}{\normalsize\bfseries}
\titleformat*{\subsubsection}{\normalsize\bfseries}
\setlength{\parskip}{0.5em plus \ww minus \ww}

\renewcommand{\quote}{\list{}{\rightmargin=\leftmargin\topsep=0pt plus \ww minus \ww}\item\relax}
%%%%%%%%%%%%%%%%%%%%%%%%%%
\usepackage{xpatch}%space equation
\xapptocmd\normalsize{%
	\abovedisplayskip=0.5em plus \ww minus \ww
	\abovedisplayshortskip=0.5em plus \ww minus \ww
	\belowdisplayskip=0.5em plus \ww minus \ww
	\belowdisplayshortskip=0.5em plus \ww minus \ww
}{}{}
%%%%%%%%%%%%%%%%%%%%%
%\setlength\intextsep{0.5em}% plus \ww minus \ww}
\setlist{topsep=0em plus \ww minus \ww, parsep=0.5em plus \ww minus \ww, itemsep=0em} 

\setcounter{topnumber}{2}
\setcounter{bottomnumber}{2}
\setcounter{totalnumber}{4}
\renewcommand{\topfraction}{0.85}
\renewcommand{\bottomfraction}{0.85}
\renewcommand{\textfraction}{0.5}
\renewcommand{\floatpagefraction}{0.7}

\renewcommand{\baselinestretch}{1.2}
%%%%%%%%%%%% flex intertext
\makeatletter
\patchcmd\set@fontsize
{\f@linespread\baselineskip}
{\f@linespread\baselineskip  plus \ww minus \ww}
{}{}
\makeatother

\usepackage{tikz}

\DeclareCaptionLabelSeparator*{spaced}{\\[1.7ex plus .01em minus .01em]}
\captionsetup[table]{textfont=bf,format=plain,labelfont=bf,justification=justified,singlelinecheck=false,labelsep=spaced,skip=0pt}
\captionsetup[figure]{labelfont=bf,justification   = RaggedRight,singlelinecheck = off,skip=5pt,labelfont=it}
%\STsetdecimalsep{{.}}\spanishdecimal{.}

%\setlength\intextsep{0.5em}% plus \ww minus \ww}
\setlength{\textfloatsep}{0.0em plus \ww minus \ww}
%https://tex.stackexchange.com/questions/39017/how-to-influence-the-position-of-float-environments-like-figure-and-table-in-lat


\usepackage[acronym,symbols,nogroupskip]{glossaries-extra}
\usepackage{glossary-superragged}
\makeglossaries

%\usepackage{dsfont}
%\usepackage{textcase}



\usepackage{textcase}
\newcommand{\autor}{Bach. Gutiérrez Brigada, AVELINO LUIS\xspace}
\newcommand{\asesor}{Lic. Canchari Soliz, Jose\xspace}

\newcommand{\expert}{\xspace}
\newcommand{\expertt}{\xspace}

\newcommand{\muestra}{30 estudiantes del colegio público \NoCaseChange{L}os \NoCaseChange{L}icenciados \NoCaseChange{A}yacucho\xspace}

\newcommand{\poblacion}{3000 estudiantes del colegio público \NoCaseChange{L}os \NoCaseChange{L}icenciados \NoCaseChange{A}yacucho\xspace}

\newcommand{\lugar}{Colegio público \NoCaseChange{L}os \NoCaseChange{L}icenciados - \NoCaseChange{A}yacucho\xspace}
%TITULO

\newcommand{\titulo}{LA MÁSCARA ANDINA PARA LA MEJORA DE LA CREATIVIDAD EN ESTUDIANTES DEL NIVEL SECUNDARIO  AYACUCHO 2023\xspace}
% VARIABLES
\newcommand{\variablei}{{{Máscara andina}}\xspace}
\newcommand{\variabled}{{{Creatividad}}\xspace} 
% DIMENSIONES
% VI
\newcommand{\dimi}{{{Dioses y deidades}}\xspace}
\newcommand{\dimii}{{{Personajes ancestrales}}\xspace}
\newcommand{\dimiii}{{{Festivas}}\xspace}
\newcommand{\dimiiii}{{{Hermanos del agua}}\xspace}
% VD 
\newcommand{\dimd}{Valoración estética\xspace}
\newcommand{\dimdd}{Pensamiento crítico\xspace}
\newcommand{\dimddd}{Sensibilidad a los problemas\xspace}
\newcommand{\dimdddd}{Elaboración\xspace}
%%%%%%%%%%%%% Metodologia
\newcommand{\enfoque}{\emph{Cuantitativa}\xspace}
\newcommand{\tipo}{\emph{Pre-experimental}\xspace}
\newcommand{\nivel}{\emph{Explicativa}\xspace}
\newcommand{\diseno}{\emph{Experimental}\xspace}

%%%%%%%%%%% INDICADORES

% INDICADORES DE LA VARIABLE INDEPENDIENTE

% Inteligencia espacial
\newcommand{\gb}{Eventos en el hogar\xspace}
\newcommand{\gbb}{Eventos en el trabajo\xspace}
\newcommand{\gbbb}{Eventos en centros de estudios\xspace}
\newcommand{\gbbbb}{Eventos en centro de esparcimiento\xspace}

\newcommand{\gbbbbb}{Personajes ficticios\xspace}
\newcommand{\gbbbbbb}{Entornos ficticios\xspace}
\newcommand{\gbbbbbbb}{Situaciones absurdas\xspace}
\newcommand{\gbbbbbbbb}{Acciones atemporales\xspace}

\newcommand{\gbbbbbbbbb}{Eventos sociales\xspace}
\newcommand{\gbbbbbbbbbb}{Personajes importantes\xspace}
\newcommand{\gbbbbbbbbbbb}{Representación de protesta\xspace}
\newcommand{\gbbbbbbbbbbbb}{Opinión publica\xspace}

\newcommand{\gbbbbbbbbbbbbb}{Temas de actualidad\xspace}
\newcommand{\gbbbbbbbbbbbbbb}{Temas de controversia\xspace}
\newcommand{\gbbbbbbbbbbbbbbb}{Escándalos sociales\xspace}
\newcommand{\gbbbbbbbbbbbbbbbb}{Corrupción social\xspace}

% INDICADORES DE LA VARIABLE DEPENDIENTE

\newcommand{\fb}{Comprende al prójimo\xspace}
\newcommand{\fbb}{Entiende emociones\xspace}
\newcommand{\fbbb}{Atiende su entorno\xspace}
\newcommand{\fbbbb}{Compasión y apoyo\xspace}

\newcommand{\fbbbbb}{Expresa pensamientos claros\xspace}
\newcommand{\fbbbbbb}{Expresa opinión responsable\xspace}
\newcommand{\fbbbbbbb}{Realiza actos autónomos\xspace}
\newcommand{\fbbbbbbbb}{Resolución de conflictos\xspace}

\newcommand{\fbbbbbbbbb}{Enfoque\xspace}
\newcommand{\fbbbbbbbbbb}{Comprensión\xspace}
\newcommand{\fbbbbbbbbbbb}{Retención\xspace}
\newcommand{\fbbbbbbbbbbbb}{Aplicación\xspace}

\newcommand{\fbbbbbbbbbbbbb}{Procesa información\xspace}
\newcommand{\fbbbbbbbbbbbbbb}{Construye ideas\xspace}
\newcommand{\fbbbbbbbbbbbbbbb}{Formula juicios\xspace}
\newcommand{\fbbbbbbbbbbbbbbbb}{Realiza reflexiones\xspace}



% PROBLEMA
\newcommand{\problema}{¿Como influye la máscara andina para la mejora de la creatividad en estudiantes del nivel secundario  de \NoCaseChange{A}yacucho 2023?\xspace}

\newcommand{\problemae}{¿Como influye la máscara andina en la  composición para la mejora de la creatividad en estudiantes del nivel secundario  de ayacucho 2023.?\xspace}

\newcommand{\problemaee}{¿Como influye la mascara andina en la originalidad para la mejora de l creatividad en estudiantes del nivel secundario de ayacucho 2023??\xspace}

\newcommand{\problemaeee}{¿Como influye la mascara andina en la expresividad para la mejora de la creatividad en estudiantes del nivel secundario de ayacucho 2023?\xspace}

\newcommand{\problemaeeee}{¿Como influye la mascara andina en el contenido para la mejora de la creatividad en estudiantes del nivel secundario de ayacucho 2023?\xspace}



% OBJETIVO
\newcommand{\objetivo}{Determinar la influencia de la máscara andina en la mejora de la creatividad en estudiantes del nivel secundario de ayacucho 2023\NoCaseChange{A}yacucho 2023\xspace}

\newcommand{\objetivoe}{Establecer la influencia de máscara andina en la composición para la mejora de la creatividad en estudiantes del nivel secundario de ayacucho 2023\xspace}

\newcommand{\objetivoee}{Demostrar la influencia de  la máscara andina en la originalidad para la mejora de la creatividad en estudiantes del nivel secundario de ayacucho 2023\xspace}

\newcommand{\objetivoeee}{Explicar la influencia de la máscara andina en la expresividad para la mejora de la creatividad en estudiantes del nivel secundario de ayacucho 2023\xspace}

\newcommand{\objetivoeeee}{Describir como influye la mascara andina en el contenido para la mejora de la creatividad en estudiantes del nivel secundario de ayacucho 2023\xspace}



% HIPOTESIS
\newcommand{\hipotesis}{La máscara andina influye en la mejora de la creatividad en estudiantes del nivel secundario de ayacucho 2023\xspace}

\newcommand{\hipotesise}{La mascara andina influye en la composición para la mejora de la creatividad en estudiantes del nivel secundario de ayacucho 2023\xspace}

\newcommand{\hipotesisee}{La mascara andina  influye  en la originalidad para la mejora de la la creatividad en estudiantes del nivel secundario 2023\xspace}

\newcommand{\hipotesiseee}{La mascara andina influye  en la expresividad para la mejora de la creatividad en estudiantes del nivel secundario de ayacucho 2023\xspace}

\newcommand{\hipotesiseeee}{La mascara andina influye  en el contenido para la mejora de la creatividad en estudiantes del nivel secundario de ayacucho 2023\xspace}



\newcommand{\N}{\mathds{N}}
\newcommand{\R}{\mathds{R}}
\newcommand{\CC}{\mathds{C}}
\newcommand{\I}{\mathds{I}}
\newcommand{\Q}{\mathds{Q}}
\newcommand{\X}{\mathds{X}}


\newcommand{\abs}[1]{\left\vert#1\right\vert}
\newcommand{\norm}[1]{\left\|#1\right\|}
\newcommand{\set}[1]{\left\{#1\right\}}
\newcommand{\seq}[1]{\left<#1\right>}
\newcommand{\co}[1]{\left[#1\right]}
\newcommand{\cc}[1]{\left(#1\right)}
\newglossaryentry{asertividad}{name=Asertividad,description={Asertividad se refiere a la capacidad de expresar tus pensamientos, necesidades y sentimientos de manera clara, honesta y respetuosa, sin dañar a los demás ni permitir que te dañen a ti mismo. La asertividad es una habilidad importante tanto en las relaciones personales como en el ámbito laboral, ya que te permite comunicarte de manera efectiva y establecer límites saludables. La comunicación asertiva implica utilizar un lenguaje claro y directo, expresar tus sentimientos de manera honesta y respetuosa, y escuchar activamente a los demás. Al ser asertivo, puedes mejorar tus relaciones, aumentar tu autoconfianza y mejorar tu bienestar emocional.}}

\newglossaryentry{color}{name=Color,description={Es una propiedad visual que se percibe en los objetos debido a la luz que reflejan o emiten. Los colores son generados por la longitud de onda de la luz visible, que se compone de los colores del espectro (rojo, naranja, amarillo, verde, azul, índigo y violeta) y sus combinaciones. El color puede ser percibido de manera subjetiva y está presente en la naturaleza, el arte, la moda, la decoración, entre otros campos.}}

\newglossaryentry{comunicacion}{name=Comunicación,description={La comunicación es un proceso de intercambio de información entre dos o más personas o entidades. Se lleva a cabo a través de diferentes medios, como lenguaje verbal, escrito, gestual, entre otros, y tiene como objetivo transmitir conocimientos, sentimientos, ideas, mensajes, y establecer un vínculo entre los interlocutores. La comunicación es esencial en todo tipo de relaciones, ya sea a nivel personal, laboral o social. Es una herramienta fundamental para solucionar problemas, tomar decisiones y establecer acuerdos. Por tanto, es importante desarrollar habilidades de comunicación efectiva para poder relacionarnos de manera satisfactoria con las personas que nos rodean. La comunicación puede ser efectiva o inefectiva, dependiendo de la forma en que se desarrolla. Para que la comunicación sea efectiva es necesario que exista una retroalimentación adecuada, una comprensión mutua, respeto y empatía entre las partes involucradas.}}

\newglossaryentry{creatividad}{name=Creatividad,description={La creatividad es la capacidad de generar ideas originales y útiles que permiten resolver problemas o satisfacer necesidades de manera innovadora. También se refiere a la habilidad de crear algo nuevo a partir de elementos ya existentes. La creatividad no se limita a las artes y la cultura, sino que se extiende a todas las áreas de la vida, incluyendo la ciencia, la tecnología, la educación, los negocios y el desarrollo humano. La creatividad puede fomentarse a través de la exploración, la experimentación, la curiosidad y el pensamiento flexible.}}

\newglossaryentry{dibujo}{name=Dibujo,description={Es una representación gráfica de una imagen o idea, realizada con distintas técnicas, como lápiz, carboncillo, tinta, acuarela o técnicas digitales. Puede ser una creación artística o técnica, y se utiliza en diferentes ámbitos, como la arquitectura, el diseño gráfico, la ilustración, la ingeniería y las artes visuales. Los dibujos pueden ser abstractos o figurativos, y pueden transmitir emociones, sensaciones y pensamientos tanto del artista como del espectador.}}

\newglossaryentry{recrear}{name=Recrear,description={Volver a crear algo que ya existió en el pasado. Imitar o representar algo o alguien de manera cercana a la realidad. Divertirse o entretenerse realizando actividades recreativas.}}

\newglossaryentry{arte}{name=Arte,description={El arte es una manifestación cultural que se expresa a través de diversas formas y disciplinas artísticas, como la pintura, la escultura, la música, la literatura, la danza, el cine, el teatro, entre otras. A través del arte, los artistas pueden expresar emociones, sentimientos, ideas y reflexiones sobre la vida y la sociedad. El arte puede ser utilizado como una herramienta de comunicación para transmitir mensajes y provocar cambios sociales, políticos o culturales. Además, el arte tiene un impacto importante en la historia y la evolución de las distintas culturas y civilizaciones. El arte también puede ser una fuente de inspiración y entretenimiento para muchas personas. Puede tener un papel importante en la educación y el desarrollo de la creatividad de los individuos, tanto en términos personales como profesionales. En resumen, el arte es una expresión cultural muy valorada en todo el mundo por su capacidad de conmover, inspirar y transformar.}}

\newglossaryentry{caricatura}{name=Caricatura,description={Las caricaturas son una forma de arte que consiste en realizar una representación exagerada y humorística de una persona o situación. Se utilizan técnicas de dibujo para resaltar ciertas características físicas o de personalidad y se suelen utilizar colores vivos y formas exageradas.}}

\newglossaryentry{comico}{name=Cómico,description={El término "cómico" se refiere a alguien o algo que hace reír o provoca risa. Puede referirse a un humorista, comediante o actor que se dedica a hacer reír a la gente a través de su actuación o sus chistes. También puede referirse a una situación o historia divertida que tiene como objetivo principal provocar risa en el público. En general, un cómico es alguien que se caracteriza por su habilidad para generar humor y alegría en los demás.}}

\newglossaryentry{esbozo}{name=Esbozo,description={En español, la palabra "esbozo" se refiere a la acción de crear una representación inicial o preliminar de algo, ya sea una idea, un dibujo o un plan. Un esbozo puede ser el punto de partida para desarrollar una obra completa, como un boceto antes de hacer un dibujo detallado o un bosquejo antes de escribir un ensayo. También se puede utilizar en sentido figurado para referirse a una descripción general o resumida de algo.}}

\newglossaryentry{periodismo}{name=Periodismo,description={El periodismo es una forma de comunicación que consiste en recopilar, verificar y difundir información sobre eventos, acontecimientos y temas de interés público. Los periodistas son responsables de investigar, escribir y presentar noticias de manera objetiva y precisa.}}

\newglossaryentry{social}{name=Social,description={El término "social" puede tener diferentes significados dependiendo del contexto en el que se utilice. En general, se refiere a todo lo relacionado con la sociedad, las relaciones entre las personas y la interacción social.}}

\newglossaryentry{interrelacion}{name=Interrelación,description={La interrelación se refiere a la conexión o relación que existe entre diferentes elementos, fenómenos o individuos. Implica que dos o más cosas están conectadas, afectándose mutuamente o influyéndose entre sí. La interrelación puede ser de diversos tipos y puede darse en distintos ámbitos, como el social, económico, político, natural, etc.}}

\newglossaryentry{conceptos}{name=Conceptos,description={Una unidad cognitiva de significado. Nace como una idea abstracta (es una construcción mental) que permite comprender las experiencias surgidas a partir de la interacción con el entorno y que, finalmente, se verbaliza (se pone en palabras).}}



\newacronym{UNSCH}{UNSCH}{Uinversidad Nacional De San Cristobal De Huamanga}
%\newacronym{IFS}{IFS}{Iterated function system (sitemas de funciones iteradas)}

\newacronym{VI}{V1}{Variable independiente}
\newacronym{VD}{V2}{Variable dependiente}
\newacronym{V}{V}{Variable}

\newacronym{D}{D}{Dimensiones}

\newacronym{M}{M}{Malo}
\newacronym{B}{B}{Bueno}
\newacronym{R}{R}{Regular}
\newacronym{E}{E}{Excelente}
\newacronym{O}{O}{Observación}

\newacronym{IT}{IT}{$i$--ésimo items de los indicadores de la variable independiente}
\newacronym{ID}{ID}{Identificador de numeración}
\newacronym{Ci}{Ci}{$i$--ésimo items de la variable independiente}
\newacronym{Pi}{Pi}{$i$--ésimo items de la variable dependiente}

%%%%%%%%%%%%%%%%%%%%%%%%%%

%\glsxtrnewsymbol[description={Promedio de las observaciones correspondientes a la dimencion $i$ de la variable dependiente, donde $j=1,2,3,4,5$, $i=1,2,3$ son los promedios en cada observacion $O^i$ de las dimension $Di$ con el método tradicional y experimental respectivamente}]{w1}{$O^i_{\overline{x}_{Dj}^T}$ y $O^i_{\overline{x}_{Dj}^E}$}

\glsxtrnewsymbol[description={Sumatoria de las varianzas de cada elemento muestral}]{var}{$\sum s^2_i$}

\glsxtrnewsymbol[description={Varianzas de cada item}]{var2}{$\sum s^2_i$}

%\glsxtrnewsymbol[description={Promedio total}]{pi1}{$\overline{x}_T$}

%\glsxtrnewsymbol[description={Calificaciones de la variable dependiente en la experimentación tradicional y experimental respectivamente}]{pi2}{$\overline{V}^T$ y $\overline{V}^E$}

%\glsxtrnewsymbol[description={Frecuencia absoluta simple}]{f}{$f$}

%\glsxtrnewsymbol[description={Promedio de los items correspodientes a la dimensión de la variable dependiente con el método tradicional y experimental respectivamente}]{xD1}{$\overline{x}_{Di}^T$ y $\overline{x}_{Di}^E$}

%\glsxtrnewsymbol[description={Promedio de los items correspodientes a la variable dependiente	con el método tradicional y experimental respectivamente}]{xD2}{$\overline{x}_{T}^T$ y $\overline{x}_{T}^E$}


\usepackage{catchfilebetweentags}
\makeatletter
\def\CatchFBT@sanitize{%
   \@sanitize
   \@makeother\{%
   \@makeother\}%
%   \endlinechar=`\^^J% <--- This line modified
}% \CatchFBT@sanitize
\makeatother



\glsaddall

%\usepackage{tocstyle}
%\usetocstyle{standard}

\usepackage{etoolbox}
\makeatletter
\pretocmd{\section}{\addtocontents{toc}{\protect\addvspace{-0.2em plus \ww minus \ww}}}{}{}
\pretocmd{\subsection}{\addtocontents{toc}{\protect\addvspace{0.5em plus \ww minus \ww}}}{}{}
\pretocmd{\subsubsection}{\addtocontents{toc}{\protect\addvspace{0.5em plus \ww minus \ww}}}{}{}
\makeatother

\usepackage{pgfplots}
\usetikzlibrary{intersections}
\pgfplotsset{compat=1.15}
\usepackage{amsmath,pgfplots,gincltex}
\begin{document}
%\SweaveOpts{concordance=TRUE}
\setlength{\abovecaptionskip}{0pt}
\setlength{\belowcaptionskip}{0pt}

%\pagenumbering{roman}
\thispagestyle{empty}

{
	\normalsize\centering\thispagestyle{empty} MINISTERIO DE EDUCACIÓN{}
	\\SUPERINTENDENCIA NACIONAL DE EDUCACIÓN SUPERIOR UNIVERSITARIA (SUNEDU)\\
	\vspace{0.65cm}
{
\large ESCUELA SUPERIOR DE FORMACIÓN ARTÍSTICA PÚBLICA\\``FELIPE GUAMÁN POMA DE AYALA''}\\
{\scriptsize CON NIVEL UNIVERSITARIO: LEY 30220}\\
\vspace{0.65cm}
\includegraphics[height=6.5cm]{logo}\\
\vspace{0.65cm}  {\large TÍTULO}\\
\vspace{0.65cm}
\MakeUppercase{\Large\titulo}\\
\vspace{0.65cm}   
{\large PRESENTADO POR\\
	\vspace{0.65cm}  
	\autor
}\\
\vspace{0.65cm}
\large TPROYECTO DE INVESTIGACIÓN PARA OPTAR EL TÍTULO PROFESIONAL DE LICENCIADA EN EDUCACIÓN ARTÍSTICA\\
\vspace{0.3cm}  ESPECIALIDAD DE ARTES PLÁSTICAS\\
\vspace{0.3cm}  AYACUCHO -- PERÚ\\
\vfill
{ 2023}

}

\newpage{}
\newpage

\flushbottom

%\renewcommand{\contentsname}{ÍNDICE DE CONTENIDOS}
%%\addcontentsline{toc}{section}{Índice de contenidos}
%\tableofcontents

%\renewcommand\listfigurename{Lista de figuras} 
%\addcontentsline{toc}{section}{Lista de figuras}
%\listoffigures

%\renewcommand\listtablename{Lista de tablas} 
%\addcontentsline{toc}{section}{Lista de tablas}
%\listoftables

%\pagenumbering{arabic}
%\setcounter{page}{1}

%\onehalfspacing
%\begin{center}
%\Large Anteproyecto de Tesis
%\end{center}
\section{TÍTULO DE LA INVESTIGACIÓN}
``\titulo''
%\tableofcontents
\section{DATOS GENERALES}
\subsection{Título} \titulo.
\subsection{Tipo de investigación} \tipo
\subsection{Autor} \autor
%\subsection{Especialidad} Matemática
\subsection{Asesor} \asesor

\section{PLANTEAMIENTO DEL PROBLEMA}
\subsection{Determinación del problema}

A nivel internacional, Uno de los problemas más comunes en la comunicación a nivel internacional es la barrera del idioma. Cuando las personas hablan diferentes idiomas, puede ser difícil entenderse mutuamente y transmitir información de manera efectiva. Esto puede conducir a malentendidos, confusiones e incluso conflictos.
Otro problema es la diferencia cultural. Cada país y cada región tiene sus propias normas, valores, creencias y formas de comunicación. Lo que puede ser considerado adecuado o normal en una cultura, puede ser ofensivo o inapropiado en otra. Esto puede generar malentendidos y hacer que las personas se sientan incomprendidas.

Además, la comunicación a nivel internacional también puede verse afectada por las diferencias en la forma de comunicarse. Algunas culturas son más directas y expresivas, mientras que otras son más indirectas y reservadas. Esto puede dificultar la interpretación de las intenciones y emociones detrás de las palabras y gestos.

También existe el problema de los estereotipos y prejuicios culturales. A menudo, las personas tienen ideas preconcebidas sobre otras culturas y países, lo que puede afectar la forma en que se comunican y se relacionan con los demás. Esto puede crear barreras en la comunicación y dificultar la construcción de relaciones sólidas y colaborativas.

En resumen, la comunicación acertiva a nivel internacional puede verse afectada por barreras del idioma, diferencias culturales, estilos de comunicación diferentes y estereotipos/prejuicios culturales. Para superar estos problemas, es importante tener una actitud abierta, flexible y respetuosa, y estar dispuesto a aprender sobre otras culturas y adaptarse a diferentes formas de comunicación.

A nivel nacional, la comunicación asertiva es un problema. Muchas personas tienen dificultades para expresar sus pensamientos y sentimientos de manera clara y respetuosa, lo que puede llevar a malentendidos y conflictos en las relaciones personales y profesionales.

Existen varios factores que contribuyen a este problema. En primer lugar, la educación en comunicación asertiva no siempre es prioritaria en el sistema escolar peruano. La mayoría de las personas aprenden a comunicarse a través de la observación y la experiencia, lo que puede resultar en patrones de comunicación poco saludables.

Además, en la cultura peruana existe una tendencia hacia la evitación del conflicto. Muchas personas prefieren evitar confrontaciones y problemas, lo que puede llevar a la falta de comunicación directa y honesta. En lugar de expresar sus necesidades y opiniones de manera clara, las personas suelen optar por la indirección y la pasividad.

Otro factor que contribuye al problema es la falta de habilidades de comunicación. Muchas personas no han sido enseñadas a escuchar activamente, a expresar sus pensamientos de manera clara y organizada, o a manejar conflictos de manera constructiva. Esto puede llevar a malentendidos, frustraciones y resentimientos en las relaciones interpersonales.

Para abordar este problema, es importante que se promueva la educación en comunicación asertiva desde la infancia. Las escuelas deben incluir programas que enseñen habilidades de comunicación, resolución de conflictos y empatía. Además, es necesario fomentar una cultura que valore la comunicación abierta, directa y respetuosa.

En resumen, la comunicación asertiva es un problema en el Perú debido a la falta de educación y habilidades en este ámbito, así como a la tendencia hacia la evitación del conflicto. Es importante promover la educación y la práctica de la comunicación asertiva para mejorar las relaciones personales y profesionales en el país.

A nivel regional, el problema de la comunicación asertiva en Ayacucho podría ser causado por varios factores, como la falta de educación en habilidades de comunicación, la barrera del idioma en comunidades indígenas, la falta de acceso a servicios de comunicación adecuados, entre otros.

En muchas áreas rurales de Ayacucho, es posible que las personas no hayan recibido una educación formal que incluya la enseñanza de habilidades de comunicación efectiva. Esto puede resultar en dificultades para expresarse claramente, entender a los demás y resolver conflictos de manera constructiva.

Además, Ayacucho es una región con una gran diversidad cultural, y muchas personas en comunidades indígenas pueden tener dificultades para comunicarse en el idioma dominante, como el español. Esto puede limitar su capacidad para expresarse y para entender y ser entendidos por los demás.
Otro problema es la falta de acceso a servicios de comunicación adecuados, como internet y telefonía celular. Esto puede dificultar la comunicación, especialmente en áreas rurales donde la infraestructura de comunicación es escasa.

Para abordar este problema, es importante proporcionar programas de educación en habilidades de comunicación en las escuelas y comunidades de Ayacucho. Además, es necesario fomentar la conservación de las lenguas indígenas y facilitar la traducción y el acceso a servicios de comunicación para todas las personas, incluso en áreas rurales\cite{gray_caricaturas_2015}.

A nivel local, El problema de la comunicación asertiva en los colegios de Ayacucho puede ser causado por varios factores. Algunos de estos factores pueden incluir:

\begin{enumerate}
	\item  Falta de habilidades de comunicación: Los estudiantes pueden no haber sido enseñados o no haber practicado habilidades de comunicación efectivas, lo que puede llevar a una comunicación no asertiva.
	\item  Barreras culturales: Las diferencias culturales pueden afectar la forma en que las personas se comunican y entienden los mensajes. Estas barreras pueden dificultar la comunicación asertiva entre diferentes grupos de estudiantes.
	\item  Falta de empatía: La falta de empatía puede dificultar la comprensión y empatía hacia los demás, lo que puede llevar a una comunicación no asertiva.
	\item  Ambiente inseguro: Un ambiente escolar inseguro puede hacer que los estudiantes se sientan inseguros o temerosos de expresar sus ideas o necesidades de manera asertiva.
	\item  Falta de enseñanza en habilidades de comunicación: Los colegios pueden no dar suficiente importancia a la enseñanza de habilidades de comunicación asertiva, lo que puede resultar en una falta de conocimiento sobre cómo comunicarse de manera efectiva.

\end{enumerate}

Para abordar este problema, es importante implementar programas y actividades que fomenten la comunicación asertiva en los colegios. Esto puede incluir talleres y sesiones de capacitación en habilidades de comunicación, promover un ambiente seguro y de apoyo para la expresión de ideas y sentimientos, y fomentar la empatía y la comprensión entre los estudiantes. También es esencial involucrar a los profesores y padres en el proceso, asegurándose de que también estén familiarizados con las habilidades de comunicación asertiva y puedan apoyar a los estudiantes en su desarrollo.

Por lo tanto el objetivo de la tesis es: Determinar como la caricatura influye para fortalecer la comunicación creativa en estudiantes del nivel secundario Ayacucho 2023. Pues Las caricaturas pueden ser una forma de comunicación muy efectiva, ya que utilizan imágenes y humor para transmitir mensajes y opiniones de manera más accesible y amigable. A través de las caricaturas, se pueden abordar temas complejos o controversiales de forma sencilla y divertida, lo que facilita que el mensaje llegue a un amplio público.

Las caricaturas también pueden ser utilizadas para la crítica social o política, utilizando personajes exagerados o estereotipados para resaltar problemas o situaciones específicas. Además, las caricaturas pueden ayudar a consolidar la identidad de una marca o de un producto, ya que a través de personajes y elementos visuales reconocibles, se puede transmitir la esencia de la empresa o del mensaje que se desea comunicar.

En resumen, las caricaturas son una forma de comunicación efectiva y atractiva, capaz de transmitir mensajes complejos o controversiales de manera accesible y divertida. \cite{grose_principios_2011}:

%\begin{displayquote}
%	\emph{Las prácticas semióticas tienen como base los elementos críticos de una
%		actividad semiótica: elección del registro en el cual representar un objeto
%		matemático, capacidad de operar con las transformaciones de tratamiento
%		y de conversión. Dar un sentido coherente a cada representación semiótica
%		de un objeto matemático dado} (p.~57)
%\end{displayquote}

\subsection{Formulación del problema}
\subsubsection{Problema General}
\problema.
\subsubsection{Problemas Específicos}
\begin{itemize}
	\item \problemae.
	\item \problemaee.
	\item \problemaeee.
	\item \problemaeeee.
\end{itemize}

%\subsection{Importancia y alcances de la investigación}
%\ExecuteMetaData[contenido.tex]{importancia}
\subsection{Justificación y viabilidad}
%El principal objetivo de la  investigación es aportar al conocimiento existente en las didácticas matemáticas, esto es, sobre la relación de la \MakeTextLowercase{\variablei}, con las \MakeTextLowercase{\variabled}, cuyos resultados se podrá sistematizarse en una propuesta para ser incorporado como conocimiento a las ciencias de la educación matemática, ya que se estaría demostrando que existe relación directa y significativa entre estas dos variables de estudio.
\subsubsection{Justificación legal}
La Constitución Política del Perú establece que la educación promueve el conocimiento y la práctica de las humanidades, la ciencia, la técnica, las artes, la educación física y el deporte, preparando para la vida y el trabajo y fomentando la solidaridad. El estado tiene la obligación de contribuir al desarrollo científico y tecnológico del país. La educación ética y cívica, así como la enseñanza de la Constitución y los derechos humanos, son obligatorias en todos los niveles educativos. La educación religiosa se imparte respetando la libertad de conciencia. Los medios de comunicación social deben cooperar con el estado en materia de educación, educación moral y cultural.

El derecho a la educación está reconocido por la Constitución Política y otras normativas nacionales, lo que implica que el Estado debe garantizar igualdad de oportunidades de acceso y retención en el sistema educativo, sin discriminación, y asegurar una educación de calidad. La Ley General de Educación Nº 28044 establece que la calidad de la educación es el nivel óptimo de formación que las personas deben alcanzar para enfrentar los desafíos del desarrollo humano, demostrando y continuando aprendiendo a lo largo de la vida. Además, se menciona la importancia de la investigación e innovación educativa, siendo responsabilidad del Ministerio de Educación promoverla dentro de su ámbito de competencia.


\subsubsection{Justificación científica}
La caricatura es una forma de comunicación visual que utiliza elementos visuales como dibujos y textos para transmitir mensajes de manera humorística o satírica. Aunque puede parecer simplemente entretenimiento, la caricatura también tiene una base científica que justifica su uso como forma de comunicación.

En primer lugar, la caricatura utiliza elementos visuales que estimulan de manera efectiva el procesamiento cognitivo del cerebro. Nuestro cerebro procesa la información visual de manera más rápida y eficiente que la información verbal. Al combinar imágenes y texto, la caricatura aprovecha esta capacidad del cerebro para procesar la información de manera más efectiva, lo que facilita la transmisión del mensaje.

Además, la caricatura utiliza elementos de humor para transmitir mensajes de manera más efectiva. La risa y el humor tienen un impacto positivo en nuestro estado de ánimo y pueden mejorar la retención y comprensión de la información. Al utilizar el humor, la caricatura logra captar la atención del espectador y mantener su interés en el mensaje que se quiere transmitir.

Adicionalmente, la caricatura es una forma de comunicación que permite expresar ideas de manera indirecta o satírica. A través del uso de la exageración, la ironía o la parodia, la caricatura puede transmitir críticas sociales o políticas de manera sutil y efectiva. Esta capacidad de la caricatura de transmitir mensajes de manera indirecta es especialmente útil en situaciones en las que la comunicación directa puede ser censurada o restringida.

En resumen, la caricatura es una forma de comunicación que utiliza elementos visuales, humor y sátira para transmitir mensajes de manera efectiva. Aprovecha la capacidad cognitiva del cerebro para procesar la información visual, utiliza el humor para captar atención y mejorar la retención de la información, y permite expresar ideas de manera indirecta o satírica. Estas bases científicas justifican el uso de la caricatura como una forma válida de comunicación.

\subsubsection{Justificación social}
La caricatura es una forma de comunicación que se utiliza para transmitir mensajes de manera efectiva y humorística. Además de entretener, las caricaturas tienen una justificación social importante ya que pueden ser utilizadas para:

\begin{enumerate}
	\item  Crítica social: Las caricaturas suelen utilizar la exageración y la sátira para resaltar problemas y defectos de la sociedad. Pueden hacer críticas políticas, económicas, culturales y sociales, y de esta manera generar reflexión y conciencia en la audiencia.
	\item Denuncia de injusticias: A través de la caricatura se pueden denunciar injusticias, abusos y problemáticas sociales. Al exponer estas situaciones de manera cómica, se captura la atención del público y se fomenta el debate y la búsqueda de soluciones.
	\item  Sensibilización: La caricatura puede ser utilizada para sensibilizar sobre problemas sociales o para llamar la atención sobre ciertos temas que necesitan ser abordados. Al usar imágenes impactantes o humorísticas, se logra captar la atención del público y generar empatía hacia ciertos asuntos.
	\item  Promoción de valores: La caricatura también puede ser una herramienta para promover valores positivos y comportamientos éticos en la sociedad. Puede transmitir mensajes de respeto, tolerancia, igualdad, solidaridad, entre otros, y contribuir así a la construcción de una sociedad más justa y equitativa.
\end{enumerate}

En resumen, la justificación social de la caricatura en la comunicación radica en su capacidad para criticar, denunciar, sensibilizar y promover valores positivos en la sociedad. A través del humor y la sátira, pueden generar conciencia y promover el cambio social.

\subsubsection{Justificación metodológica}
Una justificación metodológica de la caricatura como herramienta de comunicación puede basarse en los siguientes puntos:


\begin{enumerate}
	\item  Simplificación visual: La caricatura simplifica los rasgos y características de una persona, objeto o situación, lo que permite transmitir un mensaje de manera rápida y clara. Al simplificar la imagen, la caricatura facilita la comprensión y la asimilación del mensaje por parte del espectador.
	\item  Exageración: La exageración es una técnica muy común en la caricatura, ya que permite resaltar ciertos rasgos o aspectos de una persona o situación. La exageración ayuda a llamar la atención del espectador y a destacar los elementos más importantes o relevantes del mensaje.
	\item  Humor y sátira: La caricatura busca transmitir mensajes de manera humorística o satírica, lo que ayuda a captar la atención del espectador y a generar una respuesta emocional. El humor y la sátira pueden hacer que el mensaje sea más memorable y pueden generar reflexión y debate en torno al tema tratado.
	\item  Intencionalidad y contexto cultural: La caricatura se basa en la interpretación y el entendimiento del contexto cultural en el que se encuentra. Los elementos visuales, los colores, las situaciones y los personajes utilizados en la caricatura están diseñados de manera intencionada para transmitir un mensaje específico. La caricatura también está en constante diálogo con el entorno social y político en el que se encuentra, por lo que puede ser una respuesta a situaciones actuales o a temas de actualidad.
	\item  Libertad creativa: La caricatura ofrece a los artistas una gran libertad creativa para transmitir sus ideas y mensajes. A través de la distorsión y la exageración, los caricaturistas pueden enfatizar ciertos aspectos y criticar o satirizar situaciones o personajes sin restricciones. Esta libertad permite que la caricatura sea una forma de comunicación visual muy poderosa y efectiva.
\end{enumerate}

En conclusión, la justificación metodológica de la caricatura como herramienta de comunicación radica en su capacidad para simplificar visualmente, exagerar rasgos, utilizar el humor y la sátira, establecer diálogo con el contexto cultural y ofrecer libertad creativa a los artistas. Estos elementos hacen de la caricatura una forma de comunicación visual muy efectiva para transmitir mensajes, ideas y críticas de manera rápida, clara y memorable.


\subsubsection{Justificación práctica}
Aunque pueda parecer simple o superficial, la práctica de la caricatura tiene una serie de justificaciones prácticas que la hacen importante en el ámbito de la comunicación:

\begin{enumerate}
	\item  Síntesis visual: La caricatura condensa información en imágenes, permitiendo transmitir mensajes complejos de manera rápida y efectiva. A través de la exageración de rasgos característicos de una persona o situación, se simplifican conceptos y se facilita su comprensión.
	\item  Crítica social: La caricatura es una herramienta para hacer una crítica constructiva de la realidad social y política. Mediante la sátira, se exponen problemas y se denuncian situaciones injustas o absurdas de forma humorística, lo que puede generar conciencia o reflexión en los espectadores.
	\item  Libertad de expresión: La caricatura permite expresar opiniones de manera visual, sin necesidad de recurrir al lenguaje escrito. Esto es especialmente importante en contextos en los que la libertad de expresión puede estar limitada. Las imágenes pueden ser más poderosas y menos restrictivas que las palabras en algunos casos.
	\item  Impacto emocional: Las caricaturas tienen la capacidad de generar emociones en los espectadores, ya sea a través de la risa, la indignación o la tristeza. Esta conexión emocional puede ser de gran importancia para movilizar a la gente y promover cambios sociales.
	\item  Difusión masiva: La caricatura puede ser ampliamente difundida gracias a la facilidad de compartirla en redes sociales o imprimir en periódicos y revistas. Esto la convierte en un medio de comunicación de alcance masivo, capaz de llegar a un gran número de personas en poco tiempo.
\end{enumerate}

En resumen, la caricatura es una forma de comunicación práctica y efectiva, que sintetiza información, critica la realidad, se vale de la libertad de expresión, genera emociones y puede ser difundida masivamente. Todos estos elementos la convierten en una herramienta de gran valor en la comunicación contemporánea.




%\subsection{Limitaciones de la investigación}
%\ExecuteMetaData[contenido.tex]{justificacionteo}
%\subsubsection{Limitación metodológica}
%\ExecuteMetaData[contenido.tex]{justificacionteo}
%\subsubsection{Limitación teorica}
%\ExecuteMetaData[contenido.tex]{justificacionteo}
%\subsubsection{Limitación de recurso}
%Los recursos económicos o presupuesto serán totalmente autofinanciados.
%\subsubsection{Limitación temporal}
%\ExecuteMetaData[contenido.tex]{justificacionteo}

\section{OBJECTIVOS}
Los objetivos de la investigación se refieren a las metas o propósitos que se pretenden alcanzar a través de la realización del estudio. Estos pueden variar según la naturaleza de la investigación y los intereses del investigador, pero algunos ejemplos comunes de objetivos de investigación incluyen:

\begin{enumerate}
	\item  Identificar y describir un problema o fenómeno particular.
	\item  Determinar las causas o factores que contribuyen a un problema o fenómeno.
	\item  Explorar y analizar las opiniones, actitudes o comportamientos de un grupo de personas.
	\item  Evaluar la eficacia o impacto de un programa, política o intervención.
	\item  Desarrollar y probar una teoría o hipótesis.
	\item  Generar conocimiento nuevo o contribuir a un área de estudio existente.
	\item  Comparar y contrastar diferentes enfoques o intervenciones para determinar cuál es el más efectivo.
	\item  Recopilar datos y obtener una comprensión profunda de un tema o área específica.
	\item  Proponer recomendaciones o soluciones basadas en los hallazgos de la investigación.
	\item  Contribuir al avance de la ciencia y el conocimiento en general.
\end{enumerate}

Estos son solo algunos ejemplos, y los objetivos de investigación pueden variar ampliamente dependiendo del tema y el contexto específico de la investigación. Es importante que los objetivos sean claros, específicos y alcanzables, lo que permitirá guiar el diseño, la ejecución y la interpretación de los resultados de la investigación.

\subsubsection{Objetivo general}
\objetivo.
\subsubsection{Objetivos específicos}
\begin{itemize}
	\item \objetivoe.
	\item \objetivoee.
	\item \objetivoeee.
	\item \objetivoeeee.
\end{itemize}


\section{MARCO TEÓRICO}

\subsection{Antecedentes de estudio}
\subsubsection{Internacionales}
\cite{zurita_camacho_caricatura_2020} en la tesis  ``\usebibentry{zurita_camacho_caricatura_2020}{title}'', analiza el uso estratégico y comunicativo del personaje animado de Don Evaristo durante las administraciones municipales de Rodrigo Paz y Augusto Barrera, utilizando la metodología de la comunicación estratégica propuesta por Sandra Massoni.

Se utilizará la Versión Técnica Comunicacional (VTC) para analizar la frase núcleo del problema de comunicación, así como los componentes, síntomas y consecuencias relacionados con el espacio público urbano. También se realizará una identificación y jerarquización de los actores implicados en cada componente del problema comunicacional en ambas épocas.

Se analizarán las Marcas de Racionalidad Comunicacional propuestas por ambos personajes animados desde diferentes dimensiones de la comunicación. Además, se aplicarán indicadores de cambio social para evaluar el impacto e influencia de los procesos comunicativos, teniendo en cuenta la transformación social como respuesta a las características socioculturales locales.

El estudio se basará en entrevistas a personas clave involucradas en la creación de ambas caricaturas, con el fin de obtener distintos niveles de análisis, desde el estilo de gobierno de cada administración hasta las fortalezas y debilidades de cada caricatura en términos de comunicación estratégica. Se considerará a la comunicación estratégica como un fenómeno histórico, complejo, situacional y fluido que adopta diferentes configuraciones.


\cite{segarra_morocho_caricatura_2021} en la tesis ``\usebibentry{segarra_morocho_caricatura_2021}{title}'', se estudió el sistema narrativo de la caricatura en tres de los principales medios impresos en Ecuador. Se empleó una metodología cualitativa a través del análisis de contenido de las caricaturas publicadas durante los meses de octubre, noviembre y diciembre de 2019.

Los resultados obtenidos revelaron que la caricatura es utilizada como una forma de comunicación gráfica que emite mensajes críticos sobre las situaciones sociales, políticas, económicas, etc., de la sociedad en general. Se demostró que el tema predominante de las caricaturas en los medios analizados es la política, seguido de la sociedad, y en menor medida, el deporte y otros temas.

En conclusión, la caricatura ha sido una forma de representación exagerada de los sucesos de la vida o de un personaje específico. Se ha utilizado como recurso de comunicación gráfica en los medios impresos y forma parte de los géneros de opinión de la prensa. El presente estudio contribuyó a comprender el sistema narrativo de la caricatura y los temas que predominan en los medios analizados en Ecuador.


\cite{bravo_quezada_alisis_2018} en la tesis  ``\usebibentry{bravo_quezada_alisis_2018}{title}'', estudia la tendencia política de los diarios locales a través de la caricatura como estrategia de comunicación política. Para ello, se tomaron en cuenta las caricaturas políticas de los candidatos locales que estaban en contienda por cargos públicos en la ciudad de Cuenca.

La investigación se llevó a cabo desde agosto, mes en el que inició la campaña electoral, hasta octubre, cuando se dio el cierre de la campaña en 2004. Debido a la falta de un método de análisis de caricatura existente, se creó una guía de análisis junto con un texto de apoyo que pueda servir como aporte para la Escuela de Comunicación Social de la Universidad de Cuenca.

El estudio se divide en siete capítulos. El primero aborda el lenguaje en la comunicación humana, mientras que el segundo se centra en la caricatura como estrategia de comunicación política. El tercer capítulo examina la caricatura política a lo largo de la historia, mientras que en el cuarto se realiza el análisis de las caricaturas de los dos diarios analizados.

El quinto capítulo se enfoca en la comunicación política a través de la caricatura, y el sexto considera a la caricatura como una expresión artística de comunicación política. Finalmente, el séptimo capítulo comprende el trabajo de campo, que consiste en el análisis de la caricatura en el público lector.

En las conclusiones y recomendaciones se exponen los hallazgos del estudio, así como las sugerencias para futuras investigaciones en el tema. Con esto, se espera contribuir al análisis de la tendencia política de los diarios locales a través de la caricatura como herramienta de comunicación política.



\subsubsection{Nacionales}
\cite{felix_seras_caricatura_2013} en la tesis  ``\usebibentry{felix_seras_caricatura_2013}{title}'',  la primera parte se define la caricatura política peruana como una modalidad narrativa de opinión, explicando sus características y analizando cómo se utiliza el humor y la iconografía en este tipo de relato.

En el segundo capítulo se analizan las características de la caricatura política peruana, así como las nuevas plataformas de comunicación en las que se ha desarrollado. Se realiza un recorrido histórico de la caricatura, utilizando imágenes de publicaciones y sucesos políticos peruanos para contextualizar al lector sobre el estilo iconográfico y las diferentes cargas valorativas utilizadas desde el siglo XIX hasta la actualidad.

En la tercera parte se analiza la producción de los caricaturistas peruanos Carlos "Carlín" Tovar, Eduardo "Heduardo" Rodríguez y Andrés Edery durante las elecciones municipales a la Alcaldía de Lima. Se muestra cómo su trabajo puede considerarse una modalidad narrativa de opinión dentro del periodismo.
Finalmente, en la última parte se exponen las conclusiones generales de la investigación.


\cite{vergel_loo_imagenes_2008} en la tesis ``\usebibentry{vergel_loo_imagenes_2008}{title}'', compara el uso estratégico y comunicativo del personaje animado Don Evaristo durante las administraciones municipales de Rodrigo Paz y Augusto Barrera en Quito. A través de la metodología de la comunicación estratégica propuesta por Sandra Massoni, se analiza cómo esta caricatura se utiliza como una herramienta de comunicación estratégica para promover el cambio social en el entorno urbano. Se utiliza la Versión Técnica Comunicacional para descomponer el problema de comunicación, identificando los componentes, síntomas y consecuencias relacionadas con el espacio público urbano. Se analiza también el papel de los actores implicados en cada componente del problema comunicacional en ambas épocas. Se examinan las Marcas de Racionalidad Comunicacional propuestas por ambos personajes animados, considerando las diferentes dimensiones de la comunicación. Además, se aplican indicadores de cambio social para medir el impacto e influencia de los procesos de comunicación desde una perspectiva transformadora que tenga en cuenta las particularidades socioculturales de la localidad. El estudio se apoya en entrevistas a personas clave involucradas en la creación de ambas caricaturas, proporcionando diferentes niveles de análisis, desde el estilo de gobierno de cada administración hasta las fortalezas y debilidades de cada caricatura en términos de comunicación estratégica.

\cite{munoz_caricatura_2005} en su tesis ``\usebibentry{munoz_caricatura_2005}{title}'', se centra en el dibujo caricaturesco que se reprodujo y difundió en los medios gráficos masivos. Este género se basa en la deformación de los rasgos del rostro, reducción del cuerpo y la inclusión de objetos para resaltar las características de los personajes. El dibujo privilegia la observación cuidadosa de la estampa y los detalles, ya que una sola imagen debe transmitir toda la información necesaria. Además de la calidad de las obras, el valor de estas caricaturas radica en la cantidad de información que proporcionan sobre efemérides y el contexto en el que se crearon. Este género nos introduce en la experiencia humana a través de formas artísticas, linguísticas, religiosas, etc., y en el ámbito de la cultura política. La deformación de los rasgos humanos en este tipo de arte es una cualidad que solo los dibujantes talentosos pueden lograr, ya que requiere habilidades de las artes plásticas y una capacidad estética para experimentar y juzgar el mundo cotidiano.

\subsubsection{Regionales}
\cite{berrocal_vivanco_representacion_2015} en la tesis ``\usebibentry{berrocal_vivanco_representacion_2015}{title}'', analiza los dibujos de Guamán Poma de Ayala en su memorial Nueva Crónica y Buen Gobierno para demostrar la presencia dominante de la violencia simbólica. Se utilizó un enfoque sociocultural y se aplicaron metodologías de Hermenéutica y Semiología. Se identificaron dos tipos de violencia colonial: explícita y simbólica. Se encontró que la cultura dominante española utiliza la violencia simbólica para someter y empequeñecer a los indígenas, mientras que se representa a los españoles como poderosos y dominantes. La conclusión principal es que en los dibujos se representan la pasión, muerte y resurrección de Cristo para los indígenas, mientras que los españoles se representan como triunfantes y dominantes.

\cite{yupanqui_flujos_2019} en su tesis ``\usebibentry{yupanqui_flujos_2019}{title}'', En conclusión, el estudio realizado mostró que la caricatura política fue capaz de resistir y adaptarse a la tensión generada por el conflicto armado interno en Perú. La caricatura fue utilizada como una herramienta más para transmitir información, criticar y denunciar, pero también mostró momentos de indiferencia y falta de compromiso hacia las causas estructurales del conflicto. En los años finales de la década de los ochenta, la caricatura política se volvió más lacónica y menos proactiva.




\cite{taipe_herreras_construccion_2015} en su tesis titulado ``\usebibentry{taipe_herreras_construccion_2015}{title}'', habla sobre estereotipos que presentan a la mujer peruana como objetos sexuales y perpetúan imágenes sexistas y discriminatorias. Además, se encontró que los medios impresos desempeñan un papel fundamental en la difusión y normalización de estos estereotipos al publicar y dar visibilidad a las caricaturas.

El análisis hermenéutico permitió interpretar el significado simbólico de las caricaturas, evidenciando cómo refuerzan la cosificación de la mujer y cómo se utilizan estereotipos de género para establecer roles sociales.

Por otro lado, el análisis hemerográfico reveló que estas caricaturas son publicadas con frecuencia en los diarios nacionales, lo que demuestra su presencia y alcance en la sociedad peruana. Asimismo, se encontró que estas caricaturas se alejan de los valores morales aceptados por la sociedad, perpetuando la desigualdad de género y la objetificación de la mujer.

En conclusión, este estudio muestra que las caricaturas de corte erótico presentes en los diarios nacionales contribuyen a la construcción de estereotipos de la mujer peruana, promoviendo la cosificación y la discriminación de género. Es importante generar conciencia sobre esta problemática y promover la inclusión y el respeto hacia la mujer en los medios de comunicación.

\subsection{Bases teóricas}
\subsubsection{Contexto de la investigación}
Ayacucho es una ciudad ubicada en los Andes de Perú, conocida por su historia, cultura y arte. Durante la época colonial, la ciudad fue un importante centro de producción artística, destacándose en áreas como la pintura, escultura, arquitectura y artesanía.

En Ayacucho se desarrolló la famosa escuela pictórica conocida como la "Escuela de Ayacucho", la cual produjo una gran cantidad de pinturas religiosas de estilo barroco. Estas obras se caracterizan por su gran detalle, colorido y dramatismo, y se encuentran en numerosas iglesias y conventos de la ciudad. \cite{velasco_historia_2008}

En cuanto a la escultura, Ayacucho destaca por sus imágenes religiosas, talladas en madera o piedra. Estas esculturas son consideradas verdaderas obras de arte, con una gran expresividad y detalle en sus rasgos.

En el ámbito de la arquitectura, Ayacucho posee numerosos ejemplos de edificios coloniales, destacando la Plaza de Armas de la ciudad, rodeada de hermosas casonas coloniales y la Catedral de Ayacucho, construida en el siglo XVII.

\cite{velasco_historia_2008} menciona que en la artesanía, Ayacucho es famosa por sus tejidos, cerámicas y máscaras. Los tejidos son elaborados a mano por las comunidades indígenas de la región, utilizando técnicas tradicionales y materiales como la lana de alpaca y el algodón. Las cerámicas son elaboradas con arcilla y decoradas con motivos regionales, representando escenas de la vida cotidiana o eventos históricos. Las máscaras, por su parte, son utilizadas en festividades tradicionales y representan personajes folklóricos de la región.

En resumen, Ayacucho es una ciudad que posee una rica tradición artística, en la cual se fusionan influencias indígenas y españolas. El arte en Ayacucho ofrece una mirada única a la historia y la cultura de la región, y es un importante atractivo turístico para quienes visitan la ciudad.

\subsubsection{\variablei}%Representación de objetos matemáticos

Las caricaturas son una forma de arte que consiste en realizar una representación exagerada y humorística de una persona o situación. Se utilizan técnicas de dibujo para resaltar ciertas características físicas o de personalidad y se suelen utilizar colores vivos y formas exageradas.

Existen diferentes tipos de caricaturas, como las políticas, que se enfocan en realizar críticas o retratos humorísticos de políticos o líderes importantes. También están las caricaturas de celebridades, que se centran en resaltar las características físicas o de personalidad conocidas de figuras famosas \cite{gray_caricaturas_2015}.

Las caricaturas también están presentes en los cómics y en las tiras cómicas de los periódicos, donde se utilizan para contar historias humorísticas o transmitir mensajes de una manera divertida.

En la actualidad, las caricaturas también se han vuelto populares en la animación, donde se utilizan técnicas digitales para crear personajes y situaciones cómicas.

En resumen, las caricaturas son una forma de arte humorístico que utiliza técnicas de dibujo para resaltar características físicas o de personalidad de personas o situaciones. Son utilizadas en diferentes medios, como el arte, los cómics y la animación.

\noindent\textbf{\dimi}. De acuerdo a \cite{gray_caricaturas_2015}
son caricaturas que abordan temas relacionados con hechos reales o situaciones cotidianas comunes a los lectores, ya sea de manera simple o compleja. No todas las caricaturas de realidad son críticas, algunas tienen como único propósito hacer reír a los lectores al ridiculizar situaciones comunes, sin darles una mayor trascendencia. Se tiene los indicadores.
\begin{enumerate}
	\item  Eventos en el hogar
	\item  Eventos en el trabajo
	\item  Eventos  en centros de estudios
	\item  Eventos en centro de esparcimiento
\end{enumerate}

\noindent\textbf{\dimii}.	De acuerdo a \cite{gray_caricaturas_2015} este tipo de caricaturas abordan temas que no tienen relación alguna con los hechos de la realidad, pero no están completamente desconectadas de ella. Sus paisajes pueden ser poco usuales, ubicándose en lugares reales pero extraños o en dimensiones paralelas. Sin embargo, a través de situaciones absurdas, los personajes expresan sentimientos, emociones y pensamientos propios de la vida cotidiana. De esta manera, se ridiculizan o cuestionan formas habituales de sentir o pensar en la sociedad. La presencia de estos elementos permite al lector establecer una conexión y complicidad con el mensaje del autor. Se tiene los indicadores
\begin{enumerate}
	\item  Personajes ficticios
	\item Entornos ficticios
	\item  Situaciones absurdas
	\item  Acciones atemporales
\end{enumerate}

\noindent\textbf{\dimiii}.	De acuerdo a \cite{grose_principios_2011}, la caricatura política tiene una larga tradición y ha evolucionado junto con la sociedad y la tecnología. En épocas antiguas, se realizaban a mano y se publicaban en periódicos y revistas, mientras que en la actualidad se han popularizado en las redes sociales y se realizan de forma digital.
Algunos de los temas más comunes que se abordan en la caricatura política son la corrupción, la injusticia, las desigualdades sociales y los abusos de poder. Los caricaturistas buscan generar reflexiones y críticas a través de la exageración de los rasgos físicos o emocionales de los personajes representados.
A lo largo de la historia, la caricatura política ha enfrentado censura y represión por parte de gobiernos y grupos poderosos. Sin embargo, su impacto en la sociedad ha sido innegable, ya que ha logrado transmitir mensajes de protesta y denuncia de manera efectiva y rápida.
En conclusión, la caricatura política es un medio de comunicación importante que utiliza la sátira y la exageración para abordar temas políticos y sociales. Su poder radica en su capacidad para generar críticas y reflexiones en la opinión pública, convirtiéndola en una herramienta esencial en la construcción de una sociedad informada y participativa. Se tiene los indicadores

\begin{enumerate}
	\item  Eventos sociales
	\item  Personajes importantes
	\item  Representación de protesta
	\item  Opinión publica
\end{enumerate}

\noindent\textbf{\dimiiii}. Según \cite{gray_caricaturas_2015}	el propósito de la caricatura periodística es transmitir su mensaje de manera efectiva y provocadora, utilizando elementos visuales como la exageración de rasgos físicos y la representación simbólica de personas o ideas. A través de estas representaciones caricaturescas, se busca capturar la atención del lector y generar reflexión sobre la noticia o tema en cuestión.
Las caricaturas periodísticas suelen abordar temas de actualidad y controversia, como conflictos políticos, escándalos de corrupción, decisiones gubernamentales, problemas sociales o económicos, entre otros. A través de la sátira, el dibujante puede expresar críticas y opiniones de forma no literal, aprovechando la libertad creativa que le proporciona este género.
En los medios impresos, las caricaturas periodísticas suelen aparecer en la sección de opinión, compartiendo espacio con columnas y artículos de análisis. También se pueden encontrar en revistas y publicaciones especializadas. Con el auge de las redes sociales, las caricaturas periodísticas se han adaptado a nuevos medios, difundiéndose en plataformas digitales y alcanzando a un público más amplio.
En conclusión, la caricatura periodística combina el arte visual y el periodismo para transmitir mensajes críticos de manera humorística. Su objetivo es provocar reflexión y generar opinión sobre temas de interés público. Se tiene los indicadores


\begin{enumerate}
	\item  Temas de actualidad
	\item  Temas de controversia
	\item  Escándalos sociales
	\item  Corrupción social
\end{enumerate}


\subsubsection{\variabled}. Según \cite{canas_comunicacion_2019}, la comunicación asertiva es un estilo de comunicación en el cual se expresa de manera clara, honesta y respetuosa. Se caracteriza por ser directa y enfocada en transmitir un mensaje de manera efectiva, sin herir o agredir a la otra persona.

Una comunicación asertiva implica expresar nuestras opiniones, necesidades y emociones de manera adecuada, sin ser agresivos, pasivos o manipulativos. Busca establecer una comunicación abierta y sincera, donde ambas partes se sientan escuchadas y respetadas.

Al utilizar una comunicación asertiva, podemos transmitir nuestros pensamientos y emociones de manera clara y libre de malentendidos, evitando conflictos innecesarios y fomentando la resolución pacífica de problemas.

\noindent\textbf{\dimd}. Segun \cite{lachira_estrada_comunicacion_2020} la empatía es la capacidad de ponerse en el lugar de otra persona y comprender cómo se siente. Es la habilidad de entender y responder a las emociones y los sentimientos de los demás, y se considera una habilidad social muy valiosa para crear relaciones saludables y efectivas. La empatía requiere una combinación de atención, comprensión y sensibilidad hacia los demás para poder conectarse con ellos y mostrarles compasión y apoyo.

\begin{enumerate}
	\item  Comprende al prójimo
	\item  Entiende emociones
	\item  Atiende su entorno
	\item  Compasión y apoyo
\end{enumerate}

\noindent\textbf{\dimdd}. Según \cite{canas_comunicacion_2019},	la asertividad es la capacidad de expresar nuestros pensamientos, opiniones, sentimientos y deseos de forma clara y directa, sin agredir o ser agredidos por los demás. Consiste en ser capaces de defender nuestros derechos y actuar de forma autónoma y responsable, sin caer en la sumisión o la agresividad. La asertividad es una habilidad importante en la comunicación interpersonal y en la resolución de conflictos.

\begin{enumerate}
	\item Expresa pensamientos claros
	\item Expresa opinión responsable
	\item Realiza actos autónomos
	\item Resolución de conflictos
\end{enumerate}
\noindent\textbf{\dimddd}. Según \cite{plaza_inteligencia_2015}	La atención es un proceso cognitivo que nos permite enfocar nuestra mente en un estímulo específico o en una tarea determinada, al tiempo que ignoramos o filtramos otras distracciones. Es la capacidad de concentrarse y dirigir los recursos mentales hacia un objetivo o actividad en particular.
La atención es esencial para el aprendizaje, la memoria, la toma de decisiones y el rendimiento en general. Nos permite procesar la información de manera más eficiente, mejorar la comprensión, retención y aplicación de lo que aprendemos.
Existen diferentes tipos de atención, como la atención sostenida, que es la capacidad de mantener el enfoque en una tarea durante un período prolongado de tiempo; la atención selectiva, que implica seleccionar un estímulo relevante y filtrar las distracciones; y la atención dividida, que es la habilidad para atender a varias actividades o estímulos simultáneamente.
La atención puede verse afectada por diversos factores, como el cansancio, el estrés, la falta de motivación o la presencia de distracciones externas. Es posible mejorar y entrenar la atención a través de técnicas y estrategias como la meditación, la práctica de mindfulness, la organización y planificación, entre otras.

\begin{enumerate}
	\item Enfoque
	\item Comprensión
	\item Retención
	\item Aplicación
\end{enumerate}

\noindent\textbf{\dimdddd}. De acuerdo a \cite{canas_comunicacion_2019} el pensamiento es la capacidad cognitiva y mental de los seres humanos para procesar información, construir ideas, formular juicios, realizar reflexiones y razonar. Es una actividad mental compleja que involucra diversos procesos como la percepción, la memoria, el lenguaje, la atención y la resolución de problemas. El pensamiento puede ser consciente o inconsciente, y se manifiesta a través de diferentes formas como el razonamiento lógico, la imaginación, la intuición, la creatividad y la toma de decisiones. Es la base de la actividad mental y de la construcción del conocimiento.

\begin{enumerate}
	\item Procesa información
	\item Construye ideas
	\item Formula juicios
	\item Realiza reflexiones
\end{enumerate}
\setglossarysection{subsection}\flushbottom \printglossary[numberedsection,nonumberlist,title=Definiciones  de términos básicos]








\section{HIPÓTESIS}
\subsection{Hipótesis}
\subsubsection{Hipótesis principal}
\hipotesis.
\subsubsection{Hipótesis específicas}
\begin{itemize}
	\item \hipotesise.
	\item \hipotesisee.
	\item \hipotesiseee.
	\item \hipotesiseeee.
\end{itemize}

\section{DEFINICIÓN CONCEPTUAL Y OPERACIONALIZACÍÓN DE VARIABLES }
%\subsubsection{Operacionalización de variables}
\begin{itemize}
	\item \textbf{Variable Independiente} (\variablei)
	\item \textbf{Variable Dependiente} (\variabled)
	\item \textbf{Variable Interviniente} (Capacidad artística, interrelación social)
\end{itemize}

\subsubsection{Definición conceptual de variables}

Las variables son los conceptos, atributos o propiedades de carácter social o natural, que tienden a variar, a cambiar, a modificar para asumir niveles o escalas. Se caracterizan por ser observables, medibles y cuantificables; por tal motivo asumen valores bien definidos.

\begin{itemize}
	\item \textbf{Variable Independiente (\variablei):} Una caricatura es un retrato que exagera o distorsiona la apariencia física de una o varias personas. Es en ocasiones un retrato de la sociedad reconocible, para crear un parecido fácilmente identificable y, generalmente, humorístico. También puede tratarse de alegorías.%Es la capacidad de generar imágenes \emph{dinámicas}, mental o gráficamente (sobre una superficie o una computadora) con el objetivo de estudiar un ente abstracto o figurativo lo cual se categorizó en dos dimensiones \MakeTextLowercase{\dimi} y \MakeTextLowercase{\dimii}.

	\item \textbf{Variable Dependiente (\variabled):} Nos podemos referir a cualquier tipo de comunicación, ya que el mero hecho de comunicar es un acto creativo, en tanto que, con las informaciones que tenemos añadimos algo más ya sea consciente o inconscientemente a la información que queremos transmitir.%Es el incremento de factores que identifican la eficacia del procesamiento mental matemático en un entorno social determinado, las que deben seguir una secuencia de pasos tales como la \MakeTextLowercase{\dimd}, \MakeTextLowercase{\dimdd}, \MakeTextLowercase{\dimddd} y \MakeTextLowercase{\dimdddd}.

	\item  \textbf{Variable Interviniente (Capacidad artística, interrelación social):} La capacidad artística se refiere a las habilidades, conocimientos y talentos relacionados con la creación y expresión artística. Esto incluye habilidades en áreas como la pintura, la música, la danza, la escritura, el teatro y más.

	      La interrelación social se refiere a la capacidad de interactuar y relacionarse con otras personas de manera efectiva y satisfactoria. Implica habilidades como la comunicación, la empatía, la escucha activa y la capacidad de trabajar en equipo.

	      La capacidad artística puede tener un impacto positivo en la interrelación social. Participar en actividades artísticas puede ayudar a mejorar las habilidades de comunicación y expresión, fomentar la creatividad y promover la interacción con otras personas que comparten intereses similares.

	      Además, las actividades artísticas pueden proporcionar un espacio seguro para la expresión emocional y personal, lo que a su vez puede ayudar a mejorar el bienestar emocional y la confianza en uno mismo. Esto puede facilitar la interrelación social al brindar a las personas una mayor seguridad y comodidad al interactuar con los demás.

	      La interrelación social, a su vez, puede beneficiar la capacidad artística. Al interactuar con otras personas y compartir ideas y perspectivas, los artistas pueden ampliar su horizonte artístico y encontrar nuevas inspiraciones y enfoques. Además, la retroalimentación y el apoyo de la comunidad artística pueden ayudar a desarrollar y mejorar las habilidades artísticas.

	      En resumen, la capacidad artística y la interrelación social están estrechamente relacionadas y se refuerzan mutuamente. Participar en actividades artísticas puede mejorar las habilidades de interrelación social, al tiempo que las habilidades sociales pueden beneficiar la creatividad y el desarrollo artístico.

	      % La primera influye en poca tendencia hacia la lógica y los números y desinterés hacia los factores de productividad matemática, la segunda debido a la escasa llegada hacia la información idónea y adecuada asistida.

	      \begin{displayquote}
		      \emph{A veces llamada la variable de confusión, vincula las variables independientes y dependientes. En ciertas situaciones, la relación entre una variable independiente y una variable dependiente no puede establecerse sin la intervención de otra variable. La causa, o variable independiente, tendrá el efecto supuesto solo en presencia de una variable intermedia.} Grinnell (1988, p. 203 citado en \cite{kothari_research_2004})
	      \end{displayquote}
\end{itemize}

\subsubsection{Operacionalización de variables}
\textbf{\variablei}
Para analizar las características del proceso de aplicación de la \MakeTextLowercase{\variablei}, se establecen como dimensiones la \MakeTextLowercase{\dimi} \MakeTextLowercase{\dimii} \MakeTextLowercase{\dimiii} y \MakeTextLowercase{\dimiiii}, que cuentan con un total de 16 indicadores, que permitirán determinar el objetivo general esperado por estos procesos, por lo que se aplicara módulos de clases, talleres, exposición y discusión para mejorar el \MakeTextLowercase{\variabled}. Refiérase al Cuadro \ref{www}.

\begin{table}[ht!]
	\caption{Definición operacional de la variable independiente}
	\label{www}
	\begin{tabular}{ccccc}
		\toprule%\Xhline{1pt}
		\textbf{{VD}} & \textbf{{D}}                                      & \textbf{Indicadores}             & \textbf{Escala}                                                              & \textbf{Valores}                                                            \\\hline
		\multirow{16}{*}{\rotatebox[origin=c]{90}{\makecell[{{p{5cm}}}]{\centering\variablei}}}
		              & \multirow{3}{*}{\makecell[{{p{3cm}}}]{\ce\dimi}}  & \multirow{1}{*}{\gb}             & \multirow{16}{*}{\rotatebox[origin=c]{90}{\makecell{Cuantitativa continua}}} & \multirow{16}{*}{\rotatebox[origin=c]{90}{\makecell{Valoración cuantitaitva \\0 a 20}}} \\\cline{3-3}
		              &                                                   & \multirow{1}{*}{\gbb}                                                                                                                                                                         \\\cline{3-3}
		              &                                                   & \multirow{1}{*}{\gbbb}                                                                                                                                                                        \\\cline{3-3}
		              &                                                   & \multirow{1}{*}{\gbbb}                                                                                                                                                                        \\\cline{2-3}
		              & \multirow{4}*{\makecell[{{p{3cm}}}]{\ce\dimii}}   & \multirow{1}{*}{\gbbbb}                                                                                                                                                                       \\\cline{3-3}
		              &                                                   & \multirow{1}{*}{\gbbbbb}                                                                                                                                                                      \\\cline{3-3}
		              &                                                   & \multirow{1}{*}{\gbbbbbb}                                                                                                                                                                     \\\cline{3-3}
		              &                                                   & \multirow{1}{*}{\gbbbbbbb}                                                                                                                                                                    \\\cline{2-3}
		              & \multirow{4}*{\makecell[{{p{3cm}}}]{\ce\dimiii}}  & \multirow{1}{*}{\gbbbbbbbbb}                                                                                                                                                                  \\\cline{3-3}
		              &                                                   & \multirow{1}{*}{\gbbb}                                                                                                                                                                        \\\cline{3-3}
		              &                                                   & \multirow{1}{*}{\gbbbbbbbbb}                                                                                                                                                                  \\\cline{3-3}
		              &                                                   & \multirow{1}{*}{\gbbbbbbbbbbb}                                                                                                                                                                \\\cline{2-3}
		              & \multirow{3}*{\makecell[{{p{3cm}}}]{\ce\dimiiii}} & \multirow{1}{*}{\gbbbbbbbbbbbb}                                                                                                                                                               \\\cline{3-3}
		              &                                                   & \multirow{1}{*}{\gbbb}                                                                                                                                                                        \\\cline{3-3}
		              &                                                   & \multirow{1}{*}{\gbbb}                                                                                                                                                                        \\\cline{3-3}
		              &                                                   & \multirow{1}{*}{\gbbbbbbbbbbbbb}                                                                                                                                                              \\\cline{3-3}
		\hline
	\end{tabular}
\end{table}

\vspace{-0.1cm}

\textbf{\variabled}

El presente trabajo de investigación, se refiere a la habilidad
de incrementar el \MakeTextLowercase{\variabled} lo cual involucra procesos o indicadores como la \MakeTextLowercase{\dimd}, \MakeTextLowercase{\dimdd}, \MakeTextLowercase{\dimddd} y  \MakeTextLowercase{\dimdddd}. Se recogerá los datos usando la lista de cotejo, la fichas de observación y pruebas escritas. Para determinar el desarrollo de el \MakeTextLowercase{\variabled}, que cuentan con un total de 16 indicadores.
El cumplimiento determinará el crecimiento de la  eficacia en la \MakeTextLowercase{\variabled}, refiérase al Cuadro \ref{wwwww}.

\begin{table}[ht!]
	%\renewcommand\tabcolsep{0.15cm}\renewcommand\arraystretch{0.9}
	\caption{Definición operacional de la variable dependiente}
	\label{wwwww}
	\begin{tabular}{cllcc}
		\toprule%\Xhline{1pt}
		\textbf{{Variable}} & \textbf{{Dimensiones}}                            & \textbf{Indicadores}               & \textbf{Escala}                                                              & \textbf{Valores}                                               \\
		\hline
		\multirow{16}{*}{\rotatebox[origin=c]{90}{\makecell[{{p{3cm}}}]{\ce\variabled}}}
		                    & \multirow{4}{*}{\makecell[{{p{2.5cm}}}]{\dimd}}   & \multirow{1}{*}{\fb}               & \multirow{16}{*}{\rotatebox[origin=c]{90}{\makecell{Cuantitativa continua}}} & \multirow{16}{*}{\rotatebox[origin=c]{90}{\makecell{Valoración \\cuantitaitva de 0 a 20}}}\\\cline{3-3}
		                    &                                                   & \multirow{1}{*}{\fbb}              &                                                                              &                                                                \\\cline{3-3}
		                    &                                                   & \multirow{1}{*}{\fbbb}             &                                                                              &                                                                \\\cline{3-3}
		                    &                                                   & \multirow{1}{*}{\fbbb}             &                                                                              &                                                                \\\cline{2-3}
		                    & \multirow{4}{*}{\makecell[{{p{2.5cm}}}]{\dimdd}}  & \multirow{1}{*}{\fbbbb}            &                                                                              &                                                                \\\cline{3-3}
		                    &                                                   & \multirow{1}{*}{\fbbbbb}           &                                                                              &                                                                \\\cline{3-3}
		                    &                                                   & \multirow{1}{*}{\fbbbbbb}          &                                                                              &                                                                \\\cline{3-3}
		                    &                                                   & \multirow{1}{*}{\fbbbbbbb}         &                                                                              &                                                                \\\cline{2-3}
		                    & \multirow{4}{*}{\makecell[{{p{2.5cm}}}]{\dimddd}} & \multirow{1}{*}{\fbbbbbbbb}        &                                                                              &                                                                \\\cline{3-3}
		                    &                                                   & \multirow{1}{*}{\fbbbbbbbbb}       &                                                                              &                                                                \\\cline{3-3}
		                    &                                                   & \multirow{1}{*}{\fbbbbbbbbbb}      &                                                                              &                                                                \\\cline{3-3}
		                    &                                                   & \multirow{1}{*}{\fbbbbbbbbbbb}     &                                                                              &                                                                \\\cline{2-3}
		                    & \multirow{4}{*}{\makecell[{{p{2.5cm}}}]{\dimddd}} & \multirow{1}{*}{\fbbbbbbbbbbbb}    &                                                                              &                                                                \\\cline{3-3}
		                    &                                                   & \multirow{1}{*}{\fbbbbbbbbbbbbb}   &                                                                              &                                                                \\\cline{3-3}
		                    &                                                   & \multirow{1}{*}{\fbbbbbbbbbbbbbb}  &                                                                              &                                                                \\\cline{3-3}
		                    &                                                   & \multirow{1}{*}{\fbbbbbbbbbbbbbbb} &                                                                              &                                                                \\\bottomrule%\Xhline{1pt}
	\end{tabular}
\vspace{0.5cm}
\end{table}

%\vspace{0.5cm}

\section{ASPECTO METODOLÓGICO}

\subsection{Enfoque de la investigación}
Investigación de enfoque cuantitativo debido a que se manipulan datos numéricos con el objetivo de deducir la veracidad de las hipótesis descritas en esta investigación mediante el uso adecuado de las estadísticas descriptiva e inferencial.

De acuerdo a \cite{alzina_metodologiinvestigacion_2004} La orientación cuantitativa se fundamenta en el cálculo de características de los fenómenos sociales, lo cual presupone derivar de un cuadro conceptual adecuado al problema examinado, una cadena de proposiciones que enuncien relaciones entre si las variables experimentadas de forma justificada. Este procedimiento tiende a extender y sistematizar resultados. Según \cite{hernandez_sampieri_metodologiinvestigacion_2014}.

\begin{displayquote}
	\emph{La orientación cuantitativa se constituye, de un conjunto de términos que es secuencial y demostrativo. Cada fase antecede a la subsiguiente y no es posible omitir o evitar sucesos. La disposición es inflexible. Parte de una idea que va delimitándose y, una vez definida, se proceden a generar objetivos y cuestiones de exploración, se revisa la bibliografía y se elabora un marco o una configuración teórica. De las cuestiones se forman hipótesis y establecen variables; se bosqueja un procedimiento para experimentar la cual se conoce como diseño, se calculan las variables en un determinado contexto; se exploran los datos obtenidos, recurriendo a procesos estadísticos, y se deduce una serie de conclusiones con relación a la o las hipótesis}. (p. 97)
\end{displayquote}

\subsection{Método de investigación}

La investigación es de tipo aplicada pues esta investigación busca una teoría didáctica que pueda ser aplicada en el campo educativo una vez sea probada la veracidad y objetividad de las hipótesis planteadas.

\subsection{Tipo de investigación}
La investigación es de tipo aplicada, que tiene por finalidad contribuir al conocimiento científicos de aprendizaje y herramientas pedagógicas. De acuerdo a \cite{alzina_metodologiinvestigacion_2004}.%\citeauthor{Duval} (1999, p. 15 citado en \citeauthor{ines}, \citeyear{ines}, p.~15).

\begin{displayquote}
	\emph{Es también llamada práctica, empírica, activa y se encuentra íntimamente ligada a la investigación básica, ya que depende de sus descubrimientos y aportes teóricos para poder generar beneficios y bienestar a la sociedad. Se fundamenta en la investigación teórica; su finalidad especifica es aplicar las teorías existentes a la producción de normas y procedimientos tecnológicos para controlar situaciones  o procesos de la realidad.} (p.~15)
\end{displayquote}

Al respecto \cite{kothari_research_2004}, describe la investigación aplicada, como el manejo de las sapiencias, en la experiencia cotidiana, para emplearlos, en la totalidad de los casos, en beneficio de la humanidad.

\cite{hernandez_sampieri_metodologiinvestigacion_2014}. ``Se refiere a un estudio de investigación en
el que se manipulan deliberadamente una o más variables independientes
(supuestas causas)
para analizar las consecuencias de esa manipulación sobre una o más
variables dependientes
(supuestos efectos), dentro de una situación de control para el investigador''.

%\subsubsection{Diseño de investigación}
%La investigación es de nivel explicativa experimental, porque se busca la causa y efecto en cada una de las variables de estudio, con manipulación de la variable independiente, es decir, se manipula la variable \MakeTextLowercase{\variablei}, para efectos o influencias en la \MakeTextLowercase{\variabled}.

%Por lo que \cite{khotari} explica que en una investigación experimental, de pruebas de hipótesis, cuando un grupo está expuesto a las condiciones habituales, se denomina \emph{grupo de control} (A), pero cuando lo esta en alguna condición nueva o especial, se denomina un \emph{grupo experimental } (B). Si los grupos A y B están expuestos a programas de estudios especiales, entonces ambos grupos se denominarían grupos experimentales. Es posible diseñar estudios que incluyan sólo grupos experimentales o estudios experimentales que incluyen tanto grupos experimentales como de control.

%In an experimental hypothesis-testing research when a
%group is exposed to usual conditions, it is termed a ‘control group’, but when the group is exposed to
%some novel or special condition, it is termed an ‘experimental group’. In the above illustration, the
%Group A can be called a control group and the Group B an experimental group. If both groups A and
%B are exposed to special studies programmes, then both groups would be termed ‘experimental
%groups.’ It is possible to design studies which include only experimental groups or studies which
%include both experimental and control groups.

\subsection{Diseño de investigación}
El diseño de investigación \MakeTextLowercase{\diseno}, que se estructura de acuerdo a la Tabla \ref{figg}.
%individuos de la población iguales oportunidades de ser seleccionados, porque los grupos ya están formados con los estudiantes de la serie 200 en la \lugar.

\begin{table}[ht!]
	\caption{Diseño de investigación}
	\label{figg}
	
	\begin{tabular}{cccc}
		\toprule
		Tratamiento             & Pre & Experimento & Post \\
		\midrule
		Tratamiento 1 (Clase 1) & O1  & X1          & O2    \\
		Tratamiento 2 (Clase 2) & O3  & X2          & O4     \\
		Tratamiento 3 (Clase 3) & O5  & X3          & O6    \\
		\bottomrule
	\end{tabular}
%\vspace{0.5cm}
\end{table}

%\begin{figure}[ht!]\centering
%	\begin{tikzpicture}[> = stealth,shorten > = 2pt,semithick]
%\node[] at (0,0)  (T){$M$};
%\node[] at (2,1.3)  (o1){$O_1$};
%\node[] at (5,1.3)  (w){$V_1$};
%\node[] at (2,0)  (r){$r$};
%\node[] at (4,1.3)  (o2){$O_2$};
%\node[] at (5,-1.3)  (ww){$V_2$};
%\path[->] (o1) edge node {} (o2);
%\path[->] (T) edge node {} (ww);
%\path[<-] (o1) edge node {} (w);
%\path[<-] (o2) edge node {} (ww);
%\path[->] (o1) edge node {} (r);
%\path[->] (o2) edge node {} (r);
%	\end{tikzpicture}
%	\caption{Diseño \MakeTextLowercase{\diseno}}
%	\label{figg}
%\end{figure}

Donde $M$ es la muestra, $O_1$ es la variable 1, \MakeTextLowercase{\variablei},
$O_2$ es la variable 2, \MakeTextLowercase{\variabled} y
$r$ es la relación entre variable 1 y variable 2.

El diseño de la investigación es cuasi -- experimental de un mismo grupo de trabajo con pre y pos -- prueba, porque se tomará un solo grupo experimental. De acuerdo a textcite{design} este diseño permitirá aplicar los módulos de experimentación en un periodo determinado y luego se trabajará de manera tradicional, alternando sucesivamente.

\begin{displayquote}
	\emph{Al elegir un diseño se tiene en cuenta una secuencia de procesos para aplicarlos en un determinado problema y deducir resultados que verifiquen las hipótesis del problema. En este caso se tendrá el diseño propuesto con un solo grupo, que permitirá aplicar la primera variable sobre la segunda, para fortalecer la elaboración plástica tridimensional}. \cite{kothari_research_2004}[p 243]
\end{displayquote}



\subsection{Población y muestra}

\subsubsection{Población}

Constituida por \poblacion, matriculados en el 2023. Los criterios que se considerarán en la selección de la muestra se describen en el Tabla \ref{apt7}.


\begin{table}[ht!]
	\caption{Criterio de inclusión y exclusión}\label{apt7}
	
	\begin{tabular}{lp{4cm}l}%\Xhline{2pt}
		\toprule 
		Población  & Aptos & No aptos \\\midrule
		\multirow{3}*{\makecell*[{{p{4cm}}}]{Estudiantes matriculados del semestre impar}} & \multirow{4}*{\makecell*[{{p{4cm}}}]{Estudiantes regulares y Asistentes puntuales}} & -Estudiantes repitentes\\
			&& -Estudiantes retirados                                                                               \\
			&& -Estudiantes del quinto superior                                                                     \\
			&& -Estudiantes no asistentes                                                                          \\
		\bottomrule
	\end{tabular}
\end{table}


De acuerdo a \cite{zacarias_metodologiinvestigacion_2020}, la población o universo es el agregado de todos los casos que concuerdan con determinadas descripciones o características, es decir es un conjunto de elementos que tiene al menos una misma característica que pueden ser medidos cuantitativamente o cualitativamente.

\subsubsection{Muestra}

En la presente investigación la muestra será no probabilística e intencional compuesta por un solo grupo experimental de \muestra, al cual se le aplicará el diseño \MakeTextLowercase{\diseno} de un mismo grupo con dos tipos de pruebas en series temporales equivalentes (Alternado).

Según \cite{hernandez_sampieri_metodologiinvestigacion_2014}, la muestra es una subcolección de elementos de la población seleccionada, de donde se obtendrá información para el desarrollo del estudio y del cual se hará medicines y observaciones de la variable dependiente e independiente, interrelacionada.

\subsubsection{Muestreo}

El muestreo será no probabilístico e intencional. Esta técnica consiste  en obtener una muestra donde los elementos  se recogen de manera aleatoria.

\cite{alzina_metodologiinvestigacion_2004} se refiere a la muestra no probabilística como aquella donde las unidades de la población tiene la misma posibilidad de ser elegidas y pertenecer a la muestra. Es deliberado, ya que esta práctica se opera en poblaciones indistintas. Aquí el especialista, conociendo la población y con buen juicio resuelve que características de análisis compondrá la muestra.


\subsection{Técnicas e instrumentos de recolección de datos}

\subsubsection{Técnica de recolección de datos}

\paragraph{Observación} Para el análisis de la influencia de la variable independiente sobre la variable dependiente se utilizó las técnicas de experimentación y observación, empleándose como instrumentos el plan de experimentación y la ficha de observación. Esta técnica que permitió recoger datos de la influencia de la variable independiente \MakeTextLowercase{\variablei} sobre la variable dependiente \MakeTextLowercase{\variabled}, en el proceso de experimentación.


%
Según \cite{paz_metodologiinvestigacion_2014} el método de observación es el método más utilizado, especialmente en estudios relacionados con ciencias. La observación se convierte en una herramienta científica y en el método de recolección de datos para el investigador, cuando sirve a un propósito de investigación formulado, se planifica y registra sistemáticamente y se somete a controles de validez y fiabilidad.

\paragraph{Prueba pedagógica} Técnica que permitió recoger de manera escrita del  \MakeTextLowercase{\variablei} logro del aprendizaje de la variable dependiente \MakeTextLowercase{\variabled}, en el proceso de la experimentación.

\cite{hernandez_sampieri_metodologiinvestigacion_2014} refiere: ``Es un proceso a través del cual se compara una unidad preestablecida y que la evaluación es un proceso que consiste en obtener infom1ación sistemática y objetiva acerca de un fenómeno e interpretar dicha información a fin de seleccionar entre distintas alternativas de decisión'' (p.~68).

\paragraph{Experimental} Técnica que permitió aplicar los módulos de experimentación de la variable \variablei.


{Según \cite{kothari_research_2004} el método de observación es el método más utilizado, especialmente en estudios relacionados con ciencias. La observación se convierte en una herramienta científica y en el método de recolección de datos para el investigador, cuando sirve a un propósito de investigación formulado, se planifica y registra sistemáticamente y se somete a controles de validez y fiabilidad.}

\subsubsection{Instrumento de recolección de datos}

\paragraph{Plan de experimentación.} Documento pedagógico que contiene el proceso de aplicación de las estrategias pictóricas (ver anexo).

\paragraph{La ficha de observación} En el cual se incluyen indicadores que permitieron conocer  los indicadores de la variable independiente planteadas que influirán en el desarrollo de los indicadores de la variable dependiente en la muestra. En este caso el proceso será observado a través del evento práctico teórico que hicieron los observados.



%La ficha de observación por \cite{bernal}, es la exploración visual de lo que sucede en un contexto existente, en un fenómeno explícito, especificando y estableciendo los hechos acertados de acuerdo con algún diseño previsto, manipula un herramienta o prontuario impreso, reservado a conseguir respuestas sobre el problema en estudio.
%
%
%
\paragraph{La ficha de opinión}
%
Este instrumento permitirá recoger la opinion de los estudiantes con respecto al uso de la \MakeTextLowercase{\variablei}, y averiguar el impacto de esta. De acuerdo a cite{opinion}. Una ficha de opinión, es una ficha en la que personas externas a la investigación escribe plasma lo que piensa en relación al texto o estudio que se está realizando.
%
%
%
Esta ficha puede tener: Número (es para tener orden de las fichas), Autor (la persona quien lo escribe), Tema de que se trata, Items o cuestionarios y Nota (en caso requerido).
%
%
%
%%URL del artículo: %http://www.ejemplode.com/13-ciencia/2411-ejemplo_de_ficha_de_opinion.html
%%Leer completo: ejemplos de Ficha de opinión
%%Los indicadores que conformaran la estructura de la ficha de
%%observación fueron elaborados a partir de la operacionalización
%%de la variable independiente, siendo las posibilidades de sus
%%respuestas: No, A veces y Sí.
\paragraph{Prueba escrita} Instrumento que permitirá recoger datos de la \MakeTextLowercase{\variablei}: \MakeTextLowercase{\dimd},  \MakeTextLowercase{\dimdd}, \MakeTextLowercase{\dimddd} y \MakeTextLowercase{\dimdddd}; instrumento elaborado de acuerdo a los indicadores descritos en el presente trabajo de investigación.

Las competencias y sus valoraciones cualitativas y cuantitativas se describen en la siguiente Tabla \ref{evaluacion}.

\begin{table}[ht!]
	%\centering
	\caption{Criterio de calificación}\label{evaluacion}
	\begin{tabular}{llll}%\Xhline{2pt}
		\toprule
		\bf \multirow{2}*[-.01cm]{Capacidades}                                                   & \bf Valoración   & \bf Valoración    \\
		                                                                                         & \bf  cualificada & \bf  cuantificada \\\midrule
		\multirow{5}*[-.01cm]{\makecell*[{{p{7cm}}}]{\centering  \dimd,  \dimdd, \dimddd y \dimddd}} & Excelente        & 17 -- 20          \\
		                                                                                         & Bueno            & 13 -- 16          \\
		                                                                                         & Regular          & 09 -- 12          \\
		                                                                                         & Malo             & 05 -- 08          \\
		                                                                                         & Deficiente       & 00 -- 04          \\
		\bottomrule
	\end{tabular}
\end{table}

\subsection{Tratamiento estadístico de los datos}
Con el objetivo de procesar adecuadamente los datos y sacar conclusiones confiables, se utiliza  las siguientes estadísticas.

\subsubsection{La estadística descriptiva} Esta estadística se utiliza en el estudio  del comportamiento de los datos, para inferir conclusiones relativas de nuestras variables y la relación entre ellas.

\begin{enumerate}
	\item Tendencia central (media, mediana y moda),
	\item dispersión o variabilidad (desviación estándar, la varianza y la regresión estándar),
	\item curtosis (leptocurtica, mesocúrtica y platicúrtica) y
	\item asimetría (asimétrica positiva, simétrica y asimétrica negativa).
\end{enumerate}

\subsubsection{La estadística inferencial} Esta estadística se utilizará para probar, la normalidad y la homocedasticidad  de los datos, además de las prueba de hipótesis.

Cabe mencionar que se hará uso del software R Sweave y Excel para el procesamiento de los datos.

\subsubsection{Validación de instrumentos}
Consistirá en validar el contenido del instrumento por el criterio de jueces o expertos (con maestría o doctorado), quienes verificarán y evaluarán la coherencia y secuencialidad de los instrumentos.



La opinión de los expertos consultados permitirá establecer la validez de los instrumentos que se empleará en la investigación. La validez se establecerá mediante el método del juicio de expertos.



\begin{adjustwidth}{1cm}{}
	\emph{La validez, en términos generales, se refiere al grado en que un instrumento mide
		realmente la variable que pretende medir. Por ejemplo, un instrumento válido para medir la inteligencia debe medir la inteligencia y no la memoria. Un método para medir
		el rendimiento bursátil tiene que medir precisamente esto y no la imagen de una empresa. Un ejemplo aunque muy obvio de completa invalidez sería intentar medir el peso de los objetos con una cinta métrica en lugar de con una báscula.} \cite[p.~200]{alzina_metodologiinvestigacion_2004}
\end{adjustwidth}

%\subsection{Procedimientos}
%\subsection{Estadísticos descriptivos}
%\subsection{Estadísticos inferenciales}

\subsubsection{Cofiabilidad de instrumentos}

Existen  diversos procedimientos para calcular la confiabilidad de un instrumento de medición. Todos utilizan formulas que producen coeficientes de confiabilidad. La mayoría de estos coeficientes pueden oscilar entre cero y uno donde un coeficiente cero significa nula confiabilidad y uno representa  un máximo de confiabilidad.

La confiabilidad de consistencia interna del instrumento, se determinará con la prueba piloto, en una muestra de 10 estudiantes que no serán miembros de la muestra, aplicando mitades de Coeficiente de Pearson y la corrección de Spearman Brow, la fórmula referencial es la siguiente: $$r_{xy}=\frac{n\Xsum xy-\cc{\Xsum x}\cc{\Xsum y}}{\sqrt{n\Xsum x^2-\cc{\Xsum x}^2}\sqrt{n\Xsum y^2-\cc{\Xsum y}^2}}$$

La corrección con Spearman Brow $$r=\frac{2r_{xy}}{1+r_{xy}}$$

Donde:
$n$ tamaño de muestra
$x$ es la puntuación de la primera mitad (Impares)
$y$ es la puntuación de la segunda mitad (Pares)
$r_{xy}$ es el coeficiente Pearson
$r$ es el coeficiente de Spearman Brow. Si $0\leq r\leq 0.8$ se considera no confiable y si $0.8< r \leq1$ se considera confiable.

\begin{displayquote}
	\emph{La confiabilidad de un instrumento de medición es el grado en que su aplicación
		repetida al mismo individuo u objeto genera resultados similares. Por ejemplo, si se midiera en este momento la temperatura ambiental usando un termómetro y éste indicara que hay $22^\circ C$, y un minuto más tarde se consultara otra vez y señalara $5^\circ C$, tres minutos después se observara nuevamente y éste indicara $40^\circ C$, dicho termómetro no sería confiable, ya que su aplicación repetida produce resultados distintos. Asimismo, si una prueba de inteligencia (Intelligence Quotient, IQ) se aplica hoy a un grupo de personas y da ciertos valores de inteligencia, se aplica un mes después y proporciona valores diferentes, al igual que en subsecuentes mediciones, tal prueba no sería confiable. Los resultados no son coherentes, pues no se puede confiar en ellos.
	} \cite[p.~61]{kothari_research_2004}\end{displayquote}

\begin{table}[ht!]
		\caption{Prueba de confiabilidad de la ficha de observación}\label{pret:1}
		%\vspace{0.5cm}
		\begin{tabular}{cccccccccccc}%\Xhline{2pt}
			\toprule 
			AI             & \multicolumn{10}{c}{ITEMS} & $\sum$                                                         \\\midrule
			               & IN1                        & IN2    & IN3 & IN4 & IN5 & IN6 & IN7 & IN8 & $\ldots$ & IN23 & \\\cline{2-11}
			1              &                            &        &     &     &     &     &     &     &          &      & \\
			2              &                            &        &     &     &     &     &     &     &          &      & \\
			3              &                            &        &     &     &     &     &     &     &          &      & \\
			$\vdots$       &                            &        &     &     &     &     &     &     &          &      & \\
			20             &                            &        &     &     &     &     &     &     &          &      & \\\midrule
			$\sum$         &                            &        &     &     &     &     &     &     &          &      & \\
			$\overline{x}$ &                            &        &     &     &     &     &     &     &          &      & \\
			$s_i^2$        &                            &        &     &     &     &     &     &     &          &      & \\
			$s_i$          &                            &        &     &     &     &     &     &     &          &      & \\\bottomrule
		\end{tabular}\\\vspace{0.1cm}
		{\normalsize Fuente: Datos obtenidos de la muestra}
\end{table}


%El pre test con los datos obtenidos mediante la calificación de los expertos en artes plásticas la ficha de observación generó ... de acuerdo a la Tabla ref{pret}, el instrumento (la ficha de observación)es confiable pues $0.8<\alpha<1$.

%$$\alpha=\frac{k}{k-1}\left(1-\frac{\sum s_i^2}{s_t^2}\right)$$


\section{ASPECTO  ADMINISTRATIVO}
\subsection{Recursos}

\subsubsection{Recursos humanos}
\begin{enumerate}
	\item Responsable de la investigación
	\item Profesor Asesor
	\item Estudiantes de la \lugar.
	\item Autoridades de el \lugar.
	\item Autoridades de la Universidad ESFAPA "FGPA".
\end{enumerate}

\subsubsection{Recursos materiales}
\begin{enumerate}
	\item Bibliografía especializada e Internet %\cite{geer}
	      %\item Materiales y equipos de laboratorio %\cite{ctan}
	\item Materiales de escritorio y computadora %\cite{rahtz:96-1}
	\item Otros materiales %\cite{girou:01:}
\end{enumerate}

\subsubsection{Recursos financieros}
\begin{enumerate}
	\item Auditorio de la \lugar
	\item Laptop y Data Show
\end{enumerate}

\subsection{Presupuesto}
Materiales y equipos (financiados) refiérase al Cuadro \ref{presup}.

\begin{table}[ht!]
	\caption{Presupuesto}\label{presup}
	\begin{tabular}{ccc}\toprule
		$\text N^\circ$ & Materiales y servicios              & \multicolumn1c{Gastos en S/} \\\midrule
		1               & Materiales didácticos               & 135.30                       \\
		2               & Impresión                           & 102.50                       \\
		3               & Libros                              & 300.90                       \\
		4               & Empastados de la tesis (4 unidades) & 80.10                        \\
		5               & Gastos de movilidad                 & 90.00                        \\
		6               & Otros gastos                        & 747.00                       \\\midrule
		                & TOTAL                               & 1455.8                       \\
		\bottomrule
	\end{tabular}
\end{table}

\subsection{Cronograma de actividades}
Temporalización-cronograma de acciones refiérase al Cuadro \ref{cronog}.
\begin{table}[ht!]

	\caption{Cronograma de acciones}\label{cronog}
	\begin{tabular}{cp{6cm}cccccccccc}
		\toprule
		\ce $\text N^\circ$ & Actividades                       & \multicolumn{10}{c}{2023}                                                                                                             \\
		\midrule
		                    &                                   & E                         & M         & J         & J         & A         & S         & O         & N         & D         & E         \\\cline{3-12}
		1                   & Análisis bibliográfico            & $\bullet$                 & $\bullet$ &           &           &           &           &           &           &           &           \\
		2                   & Elaboración del marco teórico     & $\bullet$                 & $\bullet$ &           &           &           &           &           &           &           &           \\
		3                   & Elaboración de los instrumentos   &                           & $\bullet$ & $\bullet$ &           &           &           &           &           &           &           \\
		4                   & Prueba de los instrumentos        &                           &           & $\bullet$ &           &           &           &           &           &           &           \\
		5                   & Aprobación del proyecto           &                           &           &           & $\bullet$ & $\bullet$ & $\bullet$ &           &           &           &           \\
		6                   & Recolección de datos              &                           &           &           &           &           &           & $\bullet$ & $\bullet$ &           &           \\
		7                   & Procesamiento y análisis de datos &                           &           &           &           &           &           &           & $\bullet$ &           &           \\
		8                   & Redacción de la tesis             &                           &           &           &           &           &           &           &           & $\bullet$ &           \\
		9                   & Presentación de la tesis          &                           &           &           &           &           &           &           &           & $\bullet$ &           \\
		10                  & Sustentación                      &                           &           &           &           &           &           &           &           &           & $\bullet$ \\
		\bottomrule
	\end{tabular}
\end{table}

%\begin{table}[ht!]
%\centering
%\begin{spreadtab}{{tabular}{clc}}\Xhline{2pt}
%@ RECURSOS & @ MATERIALES Y SERVICIOS &@\multicolumn1c{GASTOS EN S/} \\\midrule
%@\multirow{4}*[0cm]{Materiales} & @Materiales para la escultura          @& 135.30               \\
%@& @Impresiones                           & 102.50             \\
%@& @Libros                                & 300.90            \\
%@& @Libros                                & 300.90            \\\midrule
%@\multirow{4}*[0cm]{Servicios} & @Empastados de la tesis (4 unidades)   & 80.10              \\
%@ & @Gastos de movilidad                   & 90.00                \\
%@ & @Gastos de movilidad                   & 90.00                \\
%@ & @Gastos de movilidad                   & 90.00                \\
%@ & @Otros gastos                          & 747.00               \\\midrule
%&@TOTAL@                                    & sum([0,-1]:[0,-10])        \\
%\Xhline{2pt}
%\end{spreadtab}
%  \caption{Presupuesto}\label{pp}
%\end{table}
%
%\subsection{Potencial humano}
%Personas que participaron en la realización de la investigación
%
%\begin{table}[ht]
%\centering
%\caption{Potencial humano}\label{potencialhumano}
%\begin{spreadtab}{{tabular}{ccc}}\Xhline{2pt}
%@$\text N^\circ$& @POTENCIAL HUMANO & @CANTIDAD\\\midrule
%1  & @Alumnos ($1^\circ$ artistas)           & 12 \\
%2  &@Profesores de la especialidad de escultura & 3 \\
%3  & @Autodidactas       & 2 \\
%4  & @Estadistas         & 3 \\
%5  &@Físicos             & 2 \\
%6  &@Matemáticos         & 2 \\\midrule
%&@TOTAL& sum(c2:c7)    \\
%\bottomrule
%\end{spreadtab}
%\end{table}
\section{REFERENCIAS}
%\renewcommand{\bibname}{REFERENCIAS BIBLIOGRÁFICAS}
\bibliographystyle{apacite}
\bibliography{une}

\vspace{-0.3cm}

\setglossarysection{section}
\setlength{\glsdescwidth}{0.7\textwidth}
\printglossary[nopostdot,style=long,nogroupskip,nonumberlist,type=\acronymtype,title=ACRÓNIMOS]
\vspace{-0.3cm}

\setglossarysection{section}
\setlength{\glsdescwidth}{0.7\textwidth}
\printglossary[nopostdot,style=long,nogroupskip,nonumberlist,type=symbols,title=SÍMBOLOS]
\vspace{-0.3cm}

\section*{APÉNDICES}
\addcontentsline{toc}{section}{APÉNDICES}
\renewcommand{\baselinestretch}{1.0}
\pagenumbering{Roman}

\setcounter{subsection}{0}
\renewcommand{\thesubsection}{\Alph{subsection}}
%\renewcommand{\thesubsubsection}{A.\arabic{subsection}}

\numberwithin{equation}{subsection}
\numberwithin{figure}{subsection}
\numberwithin{table}{subsection}

\subsection{Ficha técnica de la investigación \label{ficha tecnica}}
\begin{enumerate}
	\item \textbf{Titulo }\titulo\textbf{ }
	\item \textbf{Año de realización }2019
	\item \textbf{Autor de la investigación }\autor
	\item \textbf{Asesor de la investigación }\asesor
	\item \textbf{Colaboradores que validaron los instrumentos }\asesor, \expert y \expertt.
	\item \textbf{Institución } \lugar
	\item \textbf{Escuela profesional }\lugar
	\item \textbf{Formato }PDF.\textbf{ }

	      %	\item \textbf{Resumen }El presente trabajo de investigación tiene como objetivo 	\MakeTextLowercase{\objetivo}. Tipo de investigación básica, nivel 	de investigación experimental explicativa, de diseño cuasiexperimental 	de un grupo con pre y posprueba en series temporales equivalentes 	(alternado); se empleo el método hipotético deductivo, experimental	y estadístico descriptivo e inferencial; el lugar de estudio fue en	la \lugar; la muestra fue no probabilística e intencional compuesta	por un solo grupo experimental de 35 estudiantes de la serie 200,	matriculados en el curso de arquitectura, del semestre impar de la	escuela profesional de Ingeniería Civil UNSCH; los datos fueron recolectados	a través de la prueba escrita y la ficha de observación; la prueba 	de validez de instrumentos se realizó a través de juicio de expertos 	y la confiabilidad, a través de prueba del Coeficiente de Pearson 	y la corrección de Spearman Brow. Se verificó la no normalidad de 	los datos, mediante la prueba de  \emph{Shapiro Wilks}; se aplicó 	la prueba de Student dos muestras relacionadas para la prueba de hipótesis, 	con un nivel de confianza del 95\% y significancia del 5\% y se concluyó, 	que \MakeTextLowercase{\hipotesis}. \\Palabras Claves: \variablei, \MakeTextLowercase{\upshape \variabled, \dimddd, \dimi y \dimii}. 
	\item \textbf{Estructura de la tesis \flushbottom }
	      \begin{enumerate}
		      \item PLANTEAMIENTO DEL PROBLEMA
		            \begin{enumerate}
			            \item Identificación y descripción del problema
			            \item Formulación del problema
			            \item Objetivos
			            \item Justificación de la investigación
		            \end{enumerate}
		      \item MARCO TEÓRICO
		            \begin{enumerate}
			            \item Antecedentes de la investigación
			            \item Bases teóricas
			            \item Definiciones de términos básicos
		            \end{enumerate}
		      \item METODOLOGÍA DE LA INVESTIGACIÓN
		            \begin{enumerate}
			            \item Sistema de hipótesis
			            \item Sistema de variables
			            \item Operacionalización de variables
			            \item Aspecto metodológico
		            \end{enumerate}
		      \item RESULTADOS DE LA INVESTIGACIÓN
		            \begin{enumerate}
			            \item Análisis e interpretación
			            \item Resultados inferenciales
			            \item Discusión de resultados
		            \end{enumerate}
		            CONCLUSIONES

		            RECOMENDACIONES

		            REFERENCIAS BIBLIOGRÁFICAS

		            ANEXOS
	      \end{enumerate}
	\item \textbf{Antecedentes }

	\item \textbf{Problema }
	      \begin{enumerate}
		      \item \textbf{Problema general.}

		            \problema.
		      \item \textbf{Problemas específicos}
		            \begin{itemize}
			            \item \problemae.
			            \item \problemaee.
			            \item \problemaeee.
			            \item \problemaeeee.
		            \end{itemize}
	      \end{enumerate}
	\item \textbf{Objetivos }
	      \begin{enumerate}
		      \item \textbf{Objetivo general}

		            \objetivo.

		      \item \textbf{Objetivos específicos }
		            \begin{itemize}
			            \item \objetivoe.
			            \item \objetivoee.
			            \item \objetivoeee.
			            \item \objetivoeeee.
		            \end{itemize}
	      \end{enumerate}
	\item \textbf{Hipótesis }
	      \begin{enumerate}
		      \item \textbf{Hipótesis general}

		            \hipotesis.
		      \item \textbf{Hipótesis específicas}
		            \begin{itemize}
			            \item \hipotesise.
			            \item \hipotesisee.
			            \item \hipotesiseee.
			            \item \hipotesiseeee.
		            \end{itemize}
	      \end{enumerate}
	\item \textbf{Orientación metodológica y diseño muestral}
	      \begin{itemize}
		      \item \textbf{Enfoque} Cuantitativo.
		      \item \textbf{Tipo} Investigación aplicada.
		      \item \textbf{Nivel de estudio} Investigación explicativa experimental.
		      \item \textbf{Diseño} Diseño \MakeTextLowercase{\diseno} de un mismo grupo
		            con pre y postprueba en series temporales equivalentes (Alternado).
		      \item \textbf{Metodología} Experimental, hipotético deductivo, estadístico
		            y analítico.
		      \item \textbf{Población} Constituida por \poblacion.
		      \item \textbf{Muestra} Constituida por \muestra.
		      \item \textbf{Muestreo} No probabilístico.
		      \item \textbf{Técnica} Observación, evaluación pedagógica y experimentación.
		      \item \textbf{Instrumentos} Ficha de observación, ficha de opinión prueba
		            escrita y módulos experimentales.
		      \item \textbf{Procesamiento} \textbf{de datos} Mediante los programas estadísticos R Sweave y Excel.
	      \end{itemize}
	\item \textbf{Análisis y presentación de resultados} Se procesará los datos
	      recogidos en el salón de clases, es decir, se sistematizará, interpretará
	      y se sintetizará para responder así a los objetivos y a la finalidad
	      de la investigación, con la perspectiva, de que la presentación de
	      los resultados de esta investigación sirva de reflexión, sobre el
	      abordaje de la \MakeTextLowercase{\variablei}, como un factor
	      fomentador del \MakeTextLowercase{\variabled}.
	\item \textbf{Conclusión }\hipotesis\textbf{.}
\end{enumerate}
\textbf{\vspace{0.5cm}
}
\begin{center}
	....................................................\\
	La investigadora\\
	\autor
\end{center}

\clearpage

\begin{landscape}
	\subsection{Matriz de consistencia\label{consistencia}}

	%\lugar -- TÍTULO: \titulo -- AUTOR: \autor

	\begin{table}[ht!]\caption{Matriz de consistencia}
		\centering \scriptsize  \renewcommand\tabcolsep{0.1cm}\renewcommand\arraystretch{1}

		\begin{tabular}{p{3.9cm}p{3.9cm}p{3.9cm}ccc} \hline \ce\bf PROBLEMAS & \ce\bf OBJETIVOS     & \ce\bf HIPÓTESIS & \ce\bf VARIABLES & \ce\bf MARCO TEÓRICO & \bf METODOLOGÍA\Tstrut \\ \hline \ce \bf PROBLEMA GENERAL &\ce \bf OBJETIVO GENERAL &\ce \bf HIPÓTESIS GENERAL &
			\multirow{8}*[0cm]{
				\begin{minipage}[t]{2.9cm}
					\begin{itemize}[itemsep=-0pt,leftmargin=*,labelsep=.02cm,] \item \textbf{Variable 1}: \variablei.
						      \begin{itemize}[itemsep=-0pt,leftmargin=*,labelsep=.02cm,topsep=-2pt]
							      \item \dimi.
							      \item \dimii
							      \item \dimiii.
							      \item \dimiiii.
						      \end{itemize} \item \textbf{Variable 2}: \variabled.
						      \begin{itemize}[itemsep=-0pt,leftmargin=*,labelsep=.02cm,topsep=-2pt]
							      \item \dimd.
							      \item \dimdd.
							      \item \dimddd.
							      \item \dimdddd.
						      \end{itemize}
					\end{itemize}
				\end{minipage}
			}                                                     & \multirow{8}*[0cm]{
				\begin{minipage}[t]{2.9cm}
					\begin{itemize}[itemsep=-0pt,leftmargin=*,labelsep=.02cm,]
						\item Contexto de la investigación.
						\item \variablei.
						\item \variabled.
						\item \dimi.
						\item \dimii.
						\item \dimiii.
						\item \dimiiii.
						\item \dimd.
						\item \dimdd.
						\item \dimddd.
						\item \dimdddd
					\end{itemize}
				\end{minipage}
			}                                                     & \multirow{22}*[0cm]{
				\begin{minipage}[t]{6cm}
					\begin{itemize}[itemsep=-0pt,leftmargin=*,labelsep=.02cm]
						\item \textbf{Enfoque de estudio:} Investigación \MakeTextLowercase{\enfoque}.
						\item \textbf{Tipo de estudio:} Investigación \MakeTextLowercase{\tipo}. \item \textbf{Nivel de estudio:} Investigación \MakeTextLowercase{\nivel}.
						\item \textbf{Diseño de estudio:} Diseño \MakeTextLowercase{\diseno} de un mismo grupo con pre y postprueba.
						\item \textbf{Metodología:} Experimental, hipotético deductivo, estadístico y analítico.
						\item \textbf{Población:} Constituida por \poblacion. \item \textbf{Muestra:} Constituida por \muestra.
						\item \textbf{Muestreo:} Probabilístico. \item \textbf{Técnica:} Observación, evaluación peda\-gógica y experimentación.
						\item \textbf{Instrumentos:} Ficha de observación, ficha de opinión prueba escrita y módulos experimentales.
						\item \textbf{Procesamiento de datos:} Mediante los programas estadísticos R Sweave y Excel.
					\end{itemize}
				\end{minipage}
			} \Tstrut                                                                                                                                                          \\ \cline{1-3} \problema&\objetivo&\hipotesis&&&\\ \cline{1-3} \ce\bf PROBLEMAS ESPECÍFICOS &\ce\bf OBJETIVOS ESPECÍFICOS &\ce\bf  HIPÓTESIS ESPECÍFICAS &&&\Tstrut\\\cline{1-3} \problemae&\objetivoe&\hipotesise&&&\\\cline{1-3}
			\problemaee                                           & \objetivoee          & \hipotesisee     &                  &                      &                        \\\cline{1-3}
			\problemaeee                                          & \objetivoeee         & \hipotesiseee    &                  &                      &                        \\\cline{1-3}
			\problemaeeee                                         & \objetivoeeee        & \hipotesiseeee   &                  &                      &                        \\ \hline
		\end{tabular}
	\end{table}
\end{landscape}



\begin{landscape}
	\subsection{Matriz de instrumentos}

	\lugar -- TÍTULO: \titulo -- AUTOR: \autor
	\begin{table}[ht!]\caption{Matriz de instrumentos de la variable \MakeTextLowercase{\variablei}}
		\centering  \renewcommand\tabcolsep{0.1cm}\renewcommand\arraystretch{1}
		\scriptsize
		\begin{tabular}{clllcc}
			\hline
			\ce\bf VARIABLE                                        & \ce\bf DIMENSIONES        & \ce\bf INDICADORES                  & \ce\bf ÍTEMS                                                                                           & \ce\bf VALORACIÓN                                                                   & \bf INSTRUMENTOS \Tstrut                                         \\\hline
			\multirow{26}{*}{\rotatebox[origin=c]{90}{\variablei}} & \multirow{6}{*}{\dimi}    & \multirow{2}{*}{\gb}                & ¿Cuáles son los eventos en el hogar más comunes en los que se realizan caricaturas?                    & \multirow{26}{*}{\rotatebox[origin=c]{90}{Excelente Bueno Regular Malo Deficiente}} & \multirow{26}{*}{\rotatebox[origin=c]{90}{Ficha de observación}} \\\cline{4-4}
			                                                       &                           &                                     & ¿Cuánto tiempo promedio se tarda en hacer una caricatura durante un evento en el hogar?                &                                                                                                                                                        \\\cline{3-4}
			                                                       &                           & \multirow{2}{*}{\gbb}               & ¿Cuáles son las principales competencias y habilidades necesarias para trabajar como caricaturista?    &                                                                                                                                                        \\\cline{4-4}
			                                                       &                           &                                     & ¿Qué herramientas y materiales son utilizados comúnmente para hacer caricaturas?                       &                                                                                                                                                        \\\cline{3-4}
			                                                       &                           & \multirow{1}{*}{\gbbb}              & ¿Cómo puedo utilizar la técnica de la caricatura para promocionar un centro de estudios?               &                                                                                                                                                        \\\cline{3-4}
			                                                       &                           & \multirow{1}{*}{\gbbbb}             & ¿Cuáles son los beneficios de realizar caricaturas como forma de esparcimiento?                        &                                                                                                                                                        \\\cline{2-4}
			                                                       & \multirow{6}{*}{\dimii}   & \multirow{2}{*}{\gbbbbb}            & ¿Cuál es la principal característica que debes tener en cuenta al realizar una caricatura de fantasía? &                                                                                                                                                        \\\cline{4-4}
			                                                       &                           &                                     & ¿Cuáles son los elementos típicos de la fantasía que se pueden incorporar en una caricatura?           &                                                                                                                                                        \\\cline{3-4}
			                                                       &                           & \multirow{1}{*}{\gbbbbbb}           & ¿Cómo puedes resaltar el aspecto fantástico de un personaje en una caricatura?                         &                                                                                                                                                        \\\cline{3-4}
			                                                       &                           & \multirow{1}{*}{\gbbbbbbb}          & ¿Cuál es la importancia de la creatividad al diseñar una caricatura de fantasía?                       &                                                                                                                                                        \\\cline{3-4}
			                                                       &                           & \multirow{2}{*}{\gbbbbbbbb}         & ¿Qué técnicas de dibujo son más efectivas para crear caricaturas de fantasía?                          &                                                                                                                                                        \\\cline{4-4}
			                                                       &                           &                                     & ¿Cuál es el impacto de las caricaturas de fantasía en el público?                                      &                                                                                                                                                        \\\cline{2-4}
			                                                       & \multirow{7}{*}{\dimiii}  & \multirow{2}{*}{\gbbbbbbbbb}        & ¿Cómo te inspiras al realizar caricaturas políticas?                                                   &                                                                                                                                                        \\\cline{4-4}
			                                                       &                           &                                     & ¿Cuál es el objetivo de realizar caricaturas políticas?                                                &                                                                                                                                                        \\\cline{3-4}
			                                                       &                           & \multirow{1}{*}{\gbbbbbbbbbb}       & ¿Qué elementos o características son más comunes en las caricaturas políticas?                         &                                                                                                                                                        \\\cline{3-4}
			                                                       &                           & \multirow{2}{*}{\gbbbbbbbbbbb}      & ¿Cuál ha sido la reacción más destacada de políticos hacia tus caricaturas?                            &                                                                                                                                                        \\\cline{4-4}
			                                                       &                           &                                     & ¿Crees que las caricaturas políticas tienen un impacto significativo en la sociedad?                   &                                                                                                                                                        \\\cline{3-4}
			                                                       &                           & \multirow{2}{*}{\gbbbbbbbbbbbb}     & ¿Hay algún límite ético al realizar caricaturas políticas?                                             &                                                                                                                                                        \\\cline{4-4}
			                                                       &                           &                                     & ¿De qué manera las caricaturas políticas pueden influir en la opinión pública?                         &                                                                                                                                                        \\\cline{2-4}
			                                                       & \multirow{7}{*}{\dimiiii} & \multirow{2}{*}{\gbbbbbbbbbbbbb}    & ¿Cuál es la importancia de incluir caricaturas en el periodismo?                                       &                                                                                                                                                        \\\cline{4-4}
			                                                       &                           &                                     & ¿Cuáles son los principales desafíos al realizar caricaturas para el periodismo?                       &                                                                                                                                                        \\\cline{3-4}
			                                                       &                           & \multirow{1}{*}{\gbbbbbbbbbbbbbb}   & ¿Cuál es el impacto que pueden tener las caricaturas en la opinión pública?                            &                                                                                                                                                        \\\cline{3-4}
			                                                       &                           & \multirow{2}{*}{\gbbbbbbbbbbbbbbb}  & ¿Cuáles son las técnicas más comunes utilizadas al realizar caricaturas para el periodismo?            &                                                                                                                                                        \\\cline{4-4}
			                                                       &                           &                                     & ¿Cuáles son los límites éticos que deben considerarse al realizar caricaturas para el periodismo?      &                                                                                                                                                        \\\cline{3-4}
			                                                       &                           & \multirow{2}{*}{\gbbbbbbbbbbbbbbbb} & ¿Cómo se selecciona el tema o acontecimiento a caricaturizar en el periodismo?                         &                                                                                                                                                        \\\cline{4-4}
			                                                       &                           &                                     & ¿Cuál es el papel del caricaturista dentro del equipo periodístico?                                    &                                                                                                                                                        \\\hline
		\end{tabular}
	\end{table}
\end{landscape}

\begin{landscape}
	%\subsection{Matriz de instrumentos de la variable \MakeTextLowercase{\variabled} \label{1www}}

	\lugar -- TÍTULO: \titulo -- AUTOR: \autor
	\begin{table}[ht!]\caption{Matriz de instrumentos de la variable \MakeTextLowercase{\variabled}}
		\centering \scriptsize\renewcommand\tabcolsep{0.1cm}\renewcommand\arraystretch{1}
		\begin{tabular}{clllcc}
			\hline
			\ce\bf VARIABLE                                        & \ce\bf DIMENSIONES        & \ce\bf INDICADORES                 & \ce\bf ÍTEMS                                                    & \ce\bf VALORACIÓN                                                                   & \bf INSTRUMENTOS \Tstrut                                         \\\hline
			\multirow{26}{*}{\rotatebox[origin=c]{90}{\variabled}} & \multirow{6}{*}{\dimd}    & \multirow{2}{*}{\fb}               & P1: ¿Utiliza variedad de colores?                               & \multirow{26}{*}{\rotatebox[origin=c]{90}{Excelente Bueno Regular Malo Deficiente}} & \multirow{26}{*}{\rotatebox[origin=c]{90}{Ficha de observación}} \\\cline{4-4}
			                                                       &                           &                                    & P2: ¿Plasma degradación entre los colores?                      &                                                                                                                                                        \\\cline{3-4}
			                                                       &                           & \multirow{2}{*}{\fbb}              & P3: ¿Utiliza formas y superficies regulares?                    &                                                                                                                                                        \\\cline{4-4}
			                                                       &                           &                                    & P4: ¿Utiliza formas irregulares someramente?                    &                                                                                                                                                        \\\cline{3-4}
			                                                       &                           & \multirow{1}{*}{\fbb}              & P3: ¿Utiliza formas y superficies regulares?                    &                                                                                                                                                        \\\cline{3-4}
			                                                       &                           & \multirow{1}{*}{\fbbb}             & P5: ¿Manifiesta degradación de texturas?                        &                                                                                                                                                        \\\cline{2-4}
			                                                       & \multirow{6}{*}{\dimdd}   & \multirow{2}{*}{\fbbbb}            & P6: ¿Distribuye adecuadamente los tonos?                        &                                                                                                                                                        \\\cline{4-4}
			                                                       &                           &                                    & P7: ¿Asocia adecuadamente los tonos?                            &                                                                                                                                                        \\\cline{3-4}
			                                                       &                           & \multirow{1}{*}{\fbbbbb}           & P8: ¿Asocia proporcionalmente los volúmenes?                    &                                                                                                                                                        \\\cline{3-4}
			                                                       &                           & \multirow{1}{*}{\fbbbbbb}          & P9: ¿Conjuga texturas diversas?                                 &                                                                                                                                                        \\\cline{3-4}
			                                                       &                           & \multirow{2}{*}{\fbbbbbbb}         & P9: ¿Conjuga texturas diversas?                                 &                                                                                                                                                        \\\cline{4-4}
			                                                       &                           &                                    & P10: ¿Homogeneiza texturas secuencialmente?                     &                                                                                                                                                        \\\cline{2-4}
			                                                       & \multirow{7}{*}{\dimddd}  & \multirow{2}{*}{\fbbbbbbbb}        & P11: ¿Representa elementos con dimensión creciente?             &                                                                                                                                                        \\\cline{4-4}
			                                                       &                           &                                    & P12: ¿Manifiesta elementos plásticos proporcionalmente?         &                                                                                                                                                        \\\cline{3-4}
			                                                       &                           & \multirow{1}{*}{\fbbbbbbbbb}       & P13: ¿Representa los elementos plásticos de manera decreciente? &                                                                                                                                                        \\\cline{3-4}
			                                                       &                           & \multirow{2}{*}{\fbbbbbbbbbb}      & P14: ¿Conjuga los elementos plásticos?                          &                                                                                                                                                        \\\cline{4-4}
			                                                       &                           &                                    & p15: ¿Distribuye elementos, constantes ligeramente?             &                                                                                                                                                        \\\cline{3-4}
			                                                       &                           & \multirow{2}{*}{\fbbbbbbbbbbb}     & P16: ¿Yuxtapone elementos plásticos mixtos?                     &                                                                                                                                                        \\\cline{4-4}
			                                                       &                           &                                    & P16: ¿Yuxtapone elementos plásticos mixtos?                     &                                                                                                                                                        \\\cline{2-4}
			                                                       & \multirow{7}{*}{\dimdddd} & \multirow{2}{*}{\fbbbbbbbbbbbb}    & P11: ¿Representa elementos con dimensión creciente?             &                                                                                                                                                        \\\cline{4-4}
			                                                       &                           &                                    & P12: ¿Manifiesta elementos plásticos proporcionalmente?         &                                                                                                                                                        \\\cline{3-4}
			                                                       &                           & \multirow{1}{*}{\fbbbbbbbbbbbbb}   & P13: ¿Representa los elementos plásticos de manera decreciente? &                                                                                                                                                        \\\cline{3-4}
			                                                       &                           & \multirow{2}{*}{\fbbbbbbbbbbbbbb}  & P14: ¿Conjuga los elementos plásticos?                          &                                                                                                                                                        \\\cline{4-4}
			                                                       &                           &                                    & p15: ¿Distribuye elementos, constantes ligeramente?             &                                                                                                                                                        \\\cline{3-4}
			                                                       &                           & \multirow{2}{*}{\fbbbbbbbbbbbbbbb} & P16: ¿Yuxtapone elementos plásticos mixtos?                     &                                                                                                                                                        \\\cline{4-4}
			                                                       &                           &                                    & P17: ¿Yuxtapone los ritmos en el conjunto?                      &                                                                                                                                                        \\\hline
		\end{tabular}
	\end{table}
\end{landscape}


\begin{landscape}
	\subsection{Matriz de instrumentos}

	\lugar -- TÍTULO: \titulo -- AUTOR: \autor
	\begin{table}[ht!]\caption{Matriz de instrumentos de la variable \MakeTextLowercase{\variablei}}
		\centering  \renewcommand\tabcolsep{0.1cm}%\renewcommand\arraystretch{1}
		\scriptsize
		\begin{tabular}{|l|l|l|l|l|l|l|l|l|l|l|l|l|l|l|l|l|l|l|l|l|l|l|l|l|l|}
			\hline
			\multirow{2}{*}{ \ce\bf ÍTEMS}                                                                         & \multicolumn{21}{c|}{ESTUDIANTES}\Tstrut                                                                                                                       \\\cline{2-22}
			                                                                                                       & E1                                       & E2 & E3 & E4 & E5 & E6 & E7 & E8 & E9 & E10 & E11 & E12 & E13 & E14 & E15 & E16 & E17 & E18 & E19 & E20 & E2\Tstrut \\\hline
			¿Cuáles son los eventos en el hogar más comunes en los que se realizan caricaturas?                    &                                          &    &    &    &    &    &    &    &    &     &     &     &     &     &     &     &     &     &     &     &           \\\hline
			¿Cuánto tiempo promedio se tarda en hacer una caricatura durante un evento en el hogar?                &                                          &    &    &    &    &    &    &    &    &     &     &     &     &     &     &     &     &     &     &     &           \\\hline
			¿Cuáles son las principales competencias y habilidades necesarias para trabajar como caricaturista?    &                                          &    &    &    &    &    &    &    &    &     &     &     &     &     &     &     &     &     &     &     &           \\\hline
			¿Qué herramientas y materiales son utilizados comúnmente para hacer caricaturas?                       &                                          &    &    &    &    &    &    &    &    &     &     &     &     &     &     &     &     &     &     &     &           \\\hline
			¿Cómo puedo utilizar la técnica de la caricatura para promocionar un centro de estudios?               &                                          &    &    &    &    &    &    &    &    &     &     &     &     &     &     &     &     &     &     &     &           \\\hline
			¿Cuáles son los beneficios de realizar caricaturas como forma de esparcimiento?                        &                                          &    &    &    &    &    &    &    &    &     &     &     &     &     &     &     &     &     &     &     &           \\\hline
			¿Cuál es la principal característica que debes tener en cuenta al realizar una caricatura de fantasía? &                                          &    &    &    &    &    &    &    &    &     &     &     &     &     &     &     &     &     &     &     &           \\\hline
			¿Cuáles son los elementos típicos de la fantasía que se pueden incorporar en una caricatura?           &                                          &    &    &    &    &    &    &    &    &     &     &     &     &     &     &     &     &     &     &     &           \\\hline
			¿Cómo puedes resaltar el aspecto fantástico de un personaje en una caricatura?                         &                                          &    &    &    &    &    &    &    &    &     &     &     &     &     &     &     &     &     &     &     &           \\\hline
			¿Cuál es la importancia de la creatividad al diseñar una caricatura de fantasía?                       &                                          &    &    &    &    &    &    &    &    &     &     &     &     &     &     &     &     &     &     &     &           \\\hline
			¿Qué técnicas de dibujo son más efectivas para crear caricaturas de fantasía?                          &                                          &    &    &    &    &    &    &    &    &     &     &     &     &     &     &     &     &     &     &     &           \\\hline
			¿Cuál es el impacto de las caricaturas de fantasía en el público?                                      &                                          &    &    &    &    &    &    &    &    &     &     &     &     &     &     &     &     &     &     &     &           \\\hline
			¿Cómo te inspiras al realizar caricaturas políticas?                                                   &                                          &    &    &    &    &    &    &    &    &     &     &     &     &     &     &     &     &     &     &     &           \\\hline
			¿Cuál es el objetivo de realizar caricaturas políticas?                                                &                                          &    &    &    &    &    &    &    &    &     &     &     &     &     &     &     &     &     &     &     &           \\\hline
			¿Qué elementos o características son más comunes en las caricaturas políticas?                         &                                          &    &    &    &    &    &    &    &    &     &     &     &     &     &     &     &     &     &     &     &           \\\hline
			¿Cuál ha sido la reacción más destacada de políticos hacia tus caricaturas?                            &                                          &    &    &    &    &    &    &    &    &     &     &     &     &     &     &     &     &     &     &     &           \\\hline
			¿Crees que las caricaturas políticas tienen un impacto significativo en la sociedad?                   &                                          &    &    &    &    &    &    &    &    &     &     &     &     &     &     &     &     &     &     &     &           \\\hline
			¿Hay algún límite ético al realizar caricaturas políticas?                                             &                                          &    &    &    &    &    &    &    &    &     &     &     &     &     &     &     &     &     &     &     &           \\\hline
			¿De qué manera las caricaturas políticas pueden influir en la opinión pública?                         &                                          &    &    &    &    &    &    &    &    &     &     &     &     &     &     &     &     &     &     &     &           \\\hline
			¿Cuál es la importancia de incluir caricaturas en el periodismo?                                       &                                          &    &    &    &    &    &    &    &    &     &     &     &     &     &     &     &     &     &     &     &           \\\hline
			¿Cuáles son los principales desafíos al realizar caricaturas para el periodismo?                       &                                          &    &    &    &    &    &    &    &    &     &     &     &     &     &     &     &     &     &     &     &           \\\hline
			¿Cuál es el impacto que pueden tener las caricaturas en la opinión pública?                            &                                          &    &    &    &    &    &    &    &    &     &     &     &     &     &     &     &     &     &     &     &           \\\hline
			¿Cuáles son las técnicas más comunes utilizadas al realizar caricaturas para el periodismo?            &                                          &    &    &    &    &    &    &    &    &     &     &     &     &     &     &     &     &     &     &     &           \\\hline
			¿Cuáles son los límites éticos que deben considerarse al realizar caricaturas para el periodismo?      &                                          &    &    &    &    &    &    &    &    &     &     &     &     &     &     &     &     &     &     &     &           \\\hline
			¿Cómo se selecciona el tema o acontecimiento a caricaturizar en el periodismo?                         &                                          &    &    &    &    &    &    &    &    &     &     &     &     &     &     &     &     &     &     &     &           \\\hline
			¿Cuál es el papel del caricaturista dentro del equipo periodístico?                                    &                                          &    &    &    &    &    &    &    &    &     &     &     &     &     &     &     &     &     &     &     &           \\\hline
		\end{tabular}
	\end{table}
\end{landscape}




\begin{landscape}
	%\subsection{Matriz de instrumentos de la variable \MakeTextLowercase{\variabled} \label{1www}}

	\lugar -- TÍTULO: \titulo -- AUTOR: \autor
	\begin{table}[ht!]\caption{Matriz de instrumentos de la variable \MakeTextLowercase{\variabled}}
		\centering \scriptsize\renewcommand\tabcolsep{0.1cm}%\renewcommand\arraystretch{1}
		\begin{tabular}{|l|l|l|l|l|l|l|l|l|l|l|l|l|l|l|l|l|l|l|l|l|l|l|l|l|l|}
			\hline
			\multirow{2}{*}{ \ce\bf ÍTEMS}                                                                         & \multicolumn{21}{c|}{ESTUDIANTES}\Tstrut                                                                                                                       \\\cline{2-22}
			                                                                                                       & E1                                       & E2 & E3 & E4 & E5 & E6 & E7 & E8 & E9 & E10 & E11 & E12 & E13 & E14 & E15 & E16 & E17 & E18 & E19 & E20 & E2\Tstrut \\\hline
			¿Qué papel juega la empatía en la comunicación interpersonal?                                          &                                          &    &    &    &    &    &    &    &    &     &     &     &     &     &     &     &     &     &     &     &           \\\hline
			¿Cómo puede la empatía mejorar la calidad de nuestras interacciones comunicativas?                     &                                          &    &    &    &    &    &    &    &    &     &     &     &     &     &     &     &     &     &     &     &           \\\hline
			¿Cuáles son algunos ejemplos de barreras para la empatía en la comunicación?                           &                                          &    &    &    &    &    &    &    &    &     &     &     &     &     &     &     &     &     &     &     &           \\\hline
			¿Qué habilidades son necesarias para practicar la empatía en la comunicación?                          &                                          &    &    &    &    &    &    &    &    &     &     &     &     &     &     &     &     &     &     &     &           \\\hline
			¿Cuál es la diferencia entre simpatía y empatía en la comunicación?                                    &                                          &    &    &    &    &    &    &    &    &     &     &     &     &     &     &     &     &     &     &     &           \\\hline
			¿Cómo podemos fomentar la empatía en nuestras interacciones cotidianas?                                &                                          &    &    &    &    &    &    &    &    &     &     &     &     &     &     &     &     &     &     &     &           \\\hline
			¿Qué es la asertividad en la comunicación?                                                             &                                          &    &    &    &    &    &    &    &    &     &     &     &     &     &     &     &     &     &     &     &           \\\hline
			¿Cuáles son los beneficios de ser asertivo en la comunicación?                                         &                                          &    &    &    &    &    &    &    &    &     &     &     &     &     &     &     &     &     &     &     &           \\\hline
			¿Cómo se diferencia la asertividad de la agresividad y la pasividad en la comunicación?                &                                          &    &    &    &    &    &    &    &    &     &     &     &     &     &     &     &     &     &     &     &           \\\hline
			¿Cuáles son algunas técnicas o estrategias para ser más asertivo en la comunicación?                   &                                          &    &    &    &    &    &    &    &    &     &     &     &     &     &     &     &     &     &     &     &           \\\hline
			¿Cuáles son los obstáculos o barreras comunes para la asertividad en la comunicación?                  &                                          &    &    &    &    &    &    &    &    &     &     &     &     &     &     &     &     &     &     &     &           \\\hline
			¿Cómo se puede desarrollar y mejorar la habilidad de ser asertivo en la comunicación?                  &                                          &    &    &    &    &    &    &    &    &     &     &     &     &     &     &     &     &     &     &     &           \\\hline
			¿Por qué es importante la atención en la comunicación?                                                 &                                          &    &    &    &    &    &    &    &    &     &     &     &     &     &     &     &     &     &     &     &           \\\hline
			¿Cuáles son las consecuencias de una atención deficiente en la comunicación?                           &                                          &    &    &    &    &    &    &    &    &     &     &     &     &     &     &     &     &     &     &     &           \\\hline
			¿Cómo podemos mejorar nuestra capacidad de atención en la comunicación?                                &                                          &    &    &    &    &    &    &    &    &     &     &     &     &     &     &     &     &     &     &     &           \\\hline
			¿Cuáles son las barreras más comunes que impiden una atención adecuada en la comunicación?             &                                          &    &    &    &    &    &    &    &    &     &     &     &     &     &     &     &     &     &     &     &           \\\hline
			¿Qué papel juega la atención en la comunicación eficaz y exitosa?                                      &                                          &    &    &    &    &    &    &    &    &     &     &     &     &     &     &     &     &     &     &     &           \\\hline
			¿Cuál es la relación entre la atención y la comprensión en la comunicación?                            &                                          &    &    &    &    &    &    &    &    &     &     &     &     &     &     &     &     &     &     &     &           \\\hline
			¿Cómo podemos cultivar una atención consciente en nuestras interacciones comunicativas?                &                                          &    &    &    &    &    &    &    &    &     &     &     &     &     &     &     &     &     &     &     &           \\\hline
			¿Cuál es el papel del pensamiento en el proceso de comunicación?                                       &                                          &    &    &    &    &    &    &    &    &     &     &     &     &     &     &     &     &     &     &     &           \\\hline
			¿Cómo influye el pensamiento en la forma en que nos expresamos y entendemos a los demás?               &                                          &    &    &    &    &    &    &    &    &     &     &     &     &     &     &     &     &     &     &     &           \\\hline
			¿Qué papel juegan los prejuicios y estereotipos en el pensamiento durante la comunicación?             &                                          &    &    &    &    &    &    &    &    &     &     &     &     &     &     &     &     &     &     &     &           \\\hline
			¿Cómo puede el pensamiento crítico mejorar la calidad de nuestra comunicación?                         &                                          &    &    &    &    &    &    &    &    &     &     &     &     &     &     &     &     &     &     &     &           \\\hline
			¿Cuál es la relación entre el pensamiento humano y la comunicación no verbal?                          &                                          &    &    &    &    &    &    &    &    &     &     &     &     &     &     &     &     &     &     &     &           \\\hline
			¿En qué medida afecta el pensamiento creativo a nuestra forma de comunicarnos?                         &                                          &    &    &    &    &    &    &    &    &     &     &     &     &     &     &     &     &     &     &     &           \\\hline
			¿Cómo puede el pensamiento reflexivo y la autorreflexión mejorar nuestras habilidades de comunicación? &                                          &    &    &    &    &    &    &    &    &     &     &     &     &     &     &     &     &     &     &     &           \\\hline
		\end{tabular}
	\end{table}
\end{landscape}

\begin{enumerate}
	\item DATOS GENERALES
	      \begin{itemize}
		      \item Apellidos y nombres del experto: \dotfill
		      \item Cargo e institución donde labora: \dotfill
		      \item Nombre de los instrumentos: Ficha de observación de la  variable indendiente (INST1) y Ficha de observación de la  variable dendiente (INST2)
		      \item Titulo de la tesis: \titulo
		      \item Autor de los instrumentos: \autor
	      \end{itemize}

	\item CRITERIOS DE VALIDACIÓN (Acuerdo: A, Descacuerdo: D)

	      \begin{center}
		      \begin{tabular}{|c|p{6.8cm}|c|c|c|c|}
			      \hline
			      \multirow{2}*[-.01cm]{INDICADORES} & \multirow{2}*[-.01cm]{CRITERIOS DE VALIDACIÓN}                                    & \multicolumn{2}{c|}{INST1} & \multicolumn{2}{c|}{INST2}         \\			      \cline{3-6}
			                                         &                                                                                   & A                          & D                          & A & D \\			      \hline
			      Claridad                           & ¿Está formulado con lenguaje claro, apropiado y cencillo?                         &                            &                            &   &   \\			      \hline
			      Objetividad                        & ¿Las preguntas realmente recojen datos de las variables y los indicadores?        &                            &                            &   &   \\			      \hline
			      Actualización                      & ¿El instrumento es adecuado para el tipo de variable en estudio?                  &                            &                            &   &   \\			      \hline
			      Organización                       & ¿La presentacion formal (tipo, tamaño, letra, etc.) del instrumento es apropiado? &                            &                            &   &   \\			      \hline
			      Suficiencia                        & ¿Las preguntas son suficientes para recoger datos de todos los indicadores?       &                            &                            &   &   \\			      \hline
			      Intencionalidad                    & ¿Los items o preguntas responden al problema y objetivos del investigación?       &                            &                            &   &   \\			      \hline
			      Consistencia                       & ¿Los items o preguntas tienen sustento teórico y científico?                      &                            &                            &   &   \\			      \hline
			      Coherencia                         & ¿Los items o preguntas tienen son comprensibles y están bien redactados?          &                            &                            &   &   \\			      \hline
			      %		Metodología	& ¿Esta formulado con lenguaje claro, apropiado y cencillo? &  &  &  &    \\
			      %		\hline
			      %		Pertinencia	& ¿Esta formulado con lenguaje claro, apropiado y cencillo? &  &  &  &    \\
			      %		\hline
		      \end{tabular}
	      \end{center}

	\item OPINION DE APLICABILIDAD

	      ----------------------------------------------------------------------------------------------------------------------------------------------------------------
\end{enumerate}

\begin{center}

	---------------------------\\
	Firma del experto\\
	Telefono: .....................\\
	Fecha: ......./......./.....	.......
\end{center}


\subsection{Plan experimental}
\subsubsection{Título} \titulo
\subsubsection{Duración} Agosto, setiembre, octubre y noviembre del 2023.
\subsubsection{Institución} \lugar
\subsubsection{Introducción}

La caricatura suele ser un retrato de una representación humorística que tiene la intención de exagerar los rasgos físicos o faciales, ropa, aspectos corporales o comportamiento del individuo con el fin de ser grotesco.  También son una herramienta de comunicación divertida que hace un buen trabajo al transmitir un mensaje, a menudo desafiando las palabras de quienes están en el poder o de control estatal. 

De hecho, es un recurso que contiene información que se puede agregar rápidamente a la conciencia general. Expone tensiones y juicios sobre personas, comunidades o eventos y, en última instancia, no solo advierte, sino que provoca un control de la realidad. En este sentido, la presente investigación es una interesante oportunidad de investigación que favorece la interpretación de símbolos en un tiempo y  espacio específicos para identificar datos que no siempre son visibles en otras fuentes de conocimiento. Sin embargo, su enfoque requiere una comprensión amplia de su contexto y desarrollo. Tiene como objetivo sistematizar compresión los antecedentes y actualidad del tema y demostrar su valor artístico, informativo y académico. 

La actual investigación titulada “La caricatura para fortalecer la comunicación creativa, Ayacucho-2023”, responde a las dificultades de la actualidad que se vienen originando y encontrando en los estudiantes del colegio San Juan ya que viendo la problemática se ha demostrado su  poca imaginación de creativa, por lo tanto, se busca mediante la caricatura fortalecer su comunicación creativa para desarrollar el potencial creativo que representan los estudiantes, se centra el aprendizaje constructivista, y los propios estudiantes están en el centro.

\subsubsection{Objetivos}

\paragraph{General}

\objetivo

\paragraph{Específicos}

\begin{enumerate}
\item \objetivoe
\item \objetivoee
\item \objetivoeee
\item \objetivoeeee
\end{enumerate}

\subsubsection{Fundamentación pedagógica}

Uno de los ejes centrales de la enseñanza de las artes plásticas en una escuela de arte, es la caricatura dibujo, base fundamental para la construcción de sus creaciones a partir de la imaginación y la creatividad, por lo tanto, los estudiantes deben tomar mayor atención a esta dimensión artística, ya que su desarrollo lleva a generar formas diversas de imágenes, lo cual corresponde el éxito profesional.

Al respecto, \cite{pestalozzi} sostiene que la educación o su método se basa en observar, medir, dibujar, escribir y así en enseñar a distinguir la forma de cada objeto.
Por otro lado, la filosofía de Montessori (XIX), tiene un carácter marcadamente constructivista. Se entiende que la persona en general aprendemos mejor mediante contacto directo, la practica y el descubrimiento que a través de la construcción directa. No obstante, algunas asignaturas, requieren clases magistrales puntuales.

\subsubsection{Fundamentación psicológica}

Al respecto se fundamenta  social de Vygotsky, en la cual señala que el resultado de los aprendizajes que acopia un estudiante se da mediante la interacción en un entorno social y la mayor influencia se da en un medio especifico, así mismo menciona que un estudiante adquiere la clara conciencia de lo que representa y sobre todo de quien es, otro aspecto favorable es el uso de los símbolos, reconocerlos los cuales contribuyen al desarrollo de un pensamiento complejo (citado en Ortiz, 2015, p. 99)

Vygotsky en sus múltiples test que elaboro para determinar aprendizajes, llega a una teoría más reconocida como la zona de desarrollo próximo, donde menciona que los estudiantes aprenden por sí mismo, y los temas específicos aprenden con la ayuda de un experto. Por lo tanto, Vygotsky quiere demostrar que un estudiante aprende en esta zona nuevas destrezas y habilidades. Por lo tanto, conducen y guían el desarrollo, van por delante lo mismo .

\subsubsection{Organización de las actividades}
Todas las actividades consideran una duracion de 90 minutos y el instrumento de evaluacion es la ficha de observación.

\begin{table}
	\caption{Organización de las actividades}
	\label{Organizacion}
	\begin{tabular}{cp{4.5cm}ccp{4cm}}\hline
	N & \bf Nombre de la actividad                                                        & \bf Fecha      & \bf Método & \bf Recursos                                  \\\hline
	1 & Introducción general: La caricatura para fortalecer la comunicación creativa. & 24/08/2023 & Activo & Aula, Fotos, Videos, Proyector  y Cuadros \\\hline
	2 & Actividad 1. Prueba de entrada, para el recojo de datos de la pre prueba. & 31/08/2003 & Activo & Taller de dibujo, Videos y fotos, Materiales para el dibujo\\\hline
	3 & Clase teórica: fundamentos teóricos de la caricatura, y elementos dl dibujo. & 7/09/2023 & Activo & Aula, Fotos, Videos, Proyector  y Cuadros \\\hline
	4 & Introducción general: La caricatura para fortalecer la comunicación creativa. & 24/08/2023 & Activo & Aula, Fotos, Videos, Proyector  y Cuadros \\\hline
	5 & Introducción general: La caricatura para fortalecer la comunicación creativa. & 24/08/2023 & Activo & Aula, Fotos, Videos, Proyector  y Cuadros \\\hline
	6 & Introducción general: La caricatura para fortalecer la comunicación creativa. & 24/08/2023 & Activo & Aula, Fotos, Videos, Proyector  y Cuadros \\\hline
	7 & Introducción general: La caricatura para fortalecer la comunicación creativa. & 24/08/2023 & Activo & Aula, Fotos, Videos, Proyector  y Cuadros \\\hline
	8 & Introducción general: La caricatura para fortalecer la comunicación creativa. & 24/08/2023 & Activo & Aula, Fotos, Videos, Proyector  y Cuadros \\\hline
    9 & Introducción general: La caricatura para fortalecer la comunicación creativa. & 24/08/2023 & Activo & Aula, Fotos, Videos, Proyector  y Cuadros \\\hline
\end{tabular}
\end{table}


\end{document}